\chapter{Complexity — Record Table}
%%% generated by the blueprint pipeline from TCSlib.Complexity.TuringMachine.Build.RecordTable
\section{Overview}

Declarations of the Lean module \texttt{TCSlib.Complexity.TuringMachine.Build.RecordTable}.

\section{Declarations}

\begin{theorem}[Length of joined nonempty fields]
\lean{Turing.joinFields_length}
\label{Turing.joinFields_length}
\leanok
\uses{Turing.fieldOffset, Turing.joinFields, Turing.length_joinFields_add_one}
\difficulty{4}
Let $\Gamma$ be a type and let $fs$ be a nonempty finite list of finite lists over
$\Gamma$. The length of the list formed by joining the fields, with one blank separator
between consecutive fields, is the sum of the field lengths plus the number of fields
minus one.
\end{theorem}

\begin{theorem}[Total field offset]
\lean{Turing.fieldOffset_length}
\label{Turing.fieldOffset_length}
\leanok
\uses{Turing.fieldOffset, Turing.joinFields, Turing.length_joinFields_add_one}
\difficulty{1}
For a nonempty finite list $fs$ of finite fields, the offset after all its fields is the
length of the joined list plus one.
\end{theorem}

\begin{definition}[Record length]
\lean{Turing.recordLength}
\label{Turing.recordLength}
\leanok
\uses{Turing.posWidth}
For natural numbers $k$, $n$, and $s$, define the record length to be
$k(2s+1)+w(n)+k+4$, where $w(n)$ denotes the position width of $n$.
\end{definition}

\begin{theorem}[Length of a joined configuration record]
\lean{Turing.joinFields_cfgRecord_length}
\label{Turing.joinFields_cfgRecord_length}
\leanok
\uses{Turing.Cfg, Turing.binDigits, Turing.cfgRecord, Turing.cfgRecord_length, Turing.joinFields, Turing.joinFields_length, Turing.recordLength}
\difficulty{3}
Let $k,s$ be natural numbers, let $Q$ and $T$ be types, let $x$ be a finite Boolean
string, let $c$ be a configuration of a $k$-tape machine on input $x$, and let $\tau\in
T$. If $w(n)$ denotes the number of binary digits of $n+1$, then the list obtained by
joining the fields of the record formed from $s$, $c$, and $\tau$ has length
$k(2s+1)+w(|x|)+k+4$.
\end{theorem}

\begin{theorem}[Nonempty configuration record fields]
\lean{Turing.cfgRecord_field_ne_nil}
\label{Turing.cfgRecord_field_ne_nil}
\leanok
\uses{Turing.Cfg, Turing.binDigits, Turing.cfgRecord, Turing.posWidth}
\difficulty{5}
For any natural number $s$, configuration $c$ of a $k$-tape machine with Boolean
work-tape symbols, and tag $\tau$, every field in the record encoding of $c$ with tag
$\tau$ is a nonempty list.
\end{theorem}

\begin{theorem}[Record of an initial configuration]
\lean{Turing.cfgRecord_init}
\label{Turing.cfgRecord_init}
\leanok
\uses{Turing.Cfg, Turing.Cfg.init, Turing.binDigits, Turing.cfgRecord, Turing.posWidth}
\difficulty{3}
For any natural numbers $k$ and $s$, Boolean input word $x$, initial state $q_0$, and
tag $\tau_0$, the record of the initial configuration consists of a state field
containing $q_0$, the binary digits of $1$ with width determined by the length of $x$,
then $k$ identical work-tape fields of length $2s+1$ whose cells are blank and whose
head mark is set exactly at offset $s$, followed by the tag field containing $\tau_0$.
\end{theorem}

\begin{theorem}[Length of a valid configuration record]
\lean{Turing.RecordValid.joinFields_length}
\label{Turing.RecordValid.joinFields_length}
\leanok
\uses{Turing.RecSym, Turing.RecordValid, Turing.joinFields, Turing.joinFields_cfgRecord_length, Turing.recordLength}
\difficulty{4}
Let $k,n,s$ be natural numbers, let $Q,T$ be types, and let $fs$ be a list of fields of
record symbols. If $fs$ is a valid record for a configuration using $k$ work tapes,
input length $n$, and window bound $s$, then joining its fields with separators produces
a list of length $k(2s+1)+\text{posWidth}(n)+k+4$.
\end{theorem}

\begin{theorem}[Characterization of valid configuration records]
\lean{Turing.recordValid_iff_exists_cfg}
\label{Turing.recordValid_iff_exists_cfg}
\leanok
\uses{Turing.Cfg, Turing.Cfg.InWindow, Turing.RecSym, Turing.RecordValid, Turing.cfgRecord}
\difficulty{4}
Fix a natural number $k$, state and tape-symbol types $Q$ and $T$, a natural number $s$,
a Boolean input word $x$, and a finite list of lists of record symbols $fs$. The record
$fs$ is valid for parameters $k$, $|x|$, and $s$ exactly when there are a configuration
$c$ on input $x$ and a tag $\tau$ such that every work-tape head of $c$ lies within
distance $s$ of the origin, every work-tape cell outside that window is blank, and the
record of $c$ with tag $\tau$ is $fs$.
\end{theorem}

\begin{definition}[Record template]
\lean{Turing.recordTemplate}
\label{Turing.recordTemplate}
\leanok
\uses{Turing.RecSym, Turing.binDigits, Turing.posWidth}
For types $Q$ and $T$ and natural numbers $n$ and $s$, define a list of two lists of
symbols: the first consists of the binary digits of $n+1$ at width
$\operatorname{posWidth}(n)$, each regarded as a digit symbol; the second consists of
$2s+1$ blank work cells, with the head mark on the cell at index $s$.
\end{definition}

\begin{theorem}[Length of a joined record template]
\lean{Turing.joinFields_recordTemplate_length}
\label{Turing.joinFields_recordTemplate_length}
\leanok
\uses{Turing.binDigits, Turing.joinFields, Turing.joinFields_length, Turing.posWidth, Turing.recordTemplate}
\difficulty{3}
Let $Q$ and $T$ be types, and let $n$ and $s$ be natural numbers. Form a record template
from the binary digits of $n+1$ in a field of width equal to the number of binary digits
of $n+1$, followed by $2s+1$ work-cell symbols; join the two fields by inserting one
separator. The resulting list has length equal to that width plus $2s+2$.
\end{theorem}

\begin{theorem}[Binary digits at positive width]
\lean{Turing.binDigits_posWidth_succ}
\label{Turing.binDigits_posWidth_succ}
\leanok
\uses{Turing.binDigits, Turing.posWidth}
\difficulty{4}
For every $n\in\mathbb{N}$, the list of bits of $n+1$ at positions from $0$ through
$w-1$, where $w$ is the length of the binary expansion of $n+1$, equals the binary digit
list of $n+1$.
\end{theorem}

\begin{definition}[Table fields]
\lean{Turing.tableFields}
\label{Turing.tableFields}
\leanok
Given a list $h$ and a list $es$ of lists of lists, define the table fields to be the
list consisting of $h$ followed by all the lists in $es$, concatenated in order.
\end{definition}

\begin{definition}[Table field offset]
\lean{Turing.tableOffset}
\label{Turing.tableOffset}
\leanok
\uses{Turing.fieldOffset, Turing.tableFields}
Given a list $h$, a list $es$ of lists of lists, and a natural number $j$, define the
table offset as the field offset in the table fields at index one plus the sum of the
lengths of the first $j$ entries of $es$.
\end{definition}

\begin{definition}[Well-formed table]
\lean{Turing.TableWF}
\label{Turing.TableWF}
\leanok
A table with header $h$, entries $es$, and parameter $m$ is well formed exactly when $h$
is nonempty and every entry in $es$ has length $m$ and consists entirely of nonempty
lists.
\end{definition}

\begin{definition}[Tape agrees with table fields]
\lean{Turing.TableAt}
\label{Turing.TableAt}
\leanok
\uses{Turing.Cfg, Turing.fieldsTape, Turing.joinFields, Turing.tableFields}
Let $d$ be a configuration of a machine with $K$ work tapes on input $x$, let $r$ be one
of its work tapes, let $o$ be an integer, and let $h$ and $es$ be lists of lists of
symbols. The contents of tape $r$ at position $o+p$ must equal the prescribed tape
encoding of the table fields at position $p$ for every integer $p$ from $-2$ through one
more than the length of the joined fields.
\end{definition}

\begin{theorem}[Nonempty fields in a well-formed table]
\lean{Turing.TableWF.field_ne_nil}
\label{Turing.TableWF.field_ne_nil}
\leanok
\uses{Turing.TableWF, Turing.tableFields}
\difficulty{4}
Let $m\in\mathbb{N}$, let $h$ be a list, and let $es$ be a list of lists of fields.
Suppose $h$ is nonempty, every entry in $es$ has length $m$, and every field in every
entry is nonempty. Then every field in the table formed by placing $h$ first and then
all entries in $es$ is nonempty.
\end{theorem}

\begin{theorem}[Offset in a table of fields]
\lean{Turing.tableOffset_eq}
\label{Turing.tableOffset_eq}
\leanok
\uses{Turing.fieldOffset, Turing.fieldOffset_cons_succ, Turing.fieldOffset_length, Turing.joinFields, Turing.tableFields, Turing.tableOffset}
\difficulty{5}
Let $\Gamma$ be a type, let $h$ be a finite list of elements of $\Gamma$, and let $es$
be a finite list of finite lists of finite lists of elements of $\Gamma$. Suppose every
entry $e$ of $es$ is nonempty. For every natural number $j$, the offset in the table
formed from header $h$ and entries $es$, through the first $j$ entries, is $|h|+1$ plus
the sum over those entries of the length of their fields joined by single separators,
plus one separator per entry.
\end{theorem}

\begin{theorem}[Length of the complete table offset]
\lean{Turing.tableOffset_length}
\label{Turing.tableOffset_length}
\leanok
\uses{Turing.fieldOffset_length, Turing.joinFields, Turing.tableFields, Turing.tableOffset}
\difficulty{3}
Let $\Gamma$ be a type, let $h$ be a finite list over $\Gamma$, and let $es$ be a finite
list of entries, each a finite list of finite lists over $\Gamma$. The offset after all
entries of the table formed from $h$ and $es$ equals the length of the joined list of
table fields plus one.
\end{theorem}

\begin{theorem}[Coverage by a joined field]
\lean{Turing.rt_field_cover}
\label{Turing.rt_field_cover}
\leanok
\uses{Turing.fieldOffset, Turing.fieldOffset_length, Turing.fieldOffset_succ, Turing.joinFields}
\difficulty{5}
Let $\Gamma$ be a type and let $fs$ be a nonempty finite list of finite lists of
elements of $\Gamma$. For every integer $p$ with $-1\le p\le$ the length of the list
formed by joining the fields with single blank separators, there is an index $j<|fs|$
such that $p$ measured relative to the starting offset of field $j$ lies between $-1$
and the length of that field, inclusive.
\end{theorem}

\begin{theorem}[Slice of a joined tape]
\lean{Turing.rt_fields_slice}
\label{Turing.rt_fields_slice}
\leanok
\uses{Turing.fieldOffset, Turing.fieldOffset_add_length_le, Turing.fieldOffset_length, Turing.fieldOffset_succ, Turing.fieldsTape, Turing.fieldsTape_field, Turing.joinFields, Turing.rt_field_cover}
Let $\Gamma$ be a type, and let $pre$, $gs$, and $post$ be lists of lists of elements
of $\Gamma$, with $gs$ nonempty. The offset of the first middle field is
$\mathrm{fieldOffset}(pre++(gs++post),|pre|)$. This offset plus the length of
$\mathrm{joinFields}(gs)$ is at most the length of $\mathrm{joinFields}(pre++(gs++post))$.
For every integer $p$ with $-1\le p\le |\mathrm{joinFields}(gs)|$, the tape content
$\mathrm{fieldsTape}(pre++(gs++post),\mathrm{fieldOffset}(pre++(gs++post),|pre|)+p)$
equals $\mathrm{fieldsTape}(gs,p)$.
\end{theorem}

\begin{theorem}[A table entry contains a record prefix]
\lean{Turing.TableAt.recordAt_take}
\label{Turing.TableAt.recordAt_take}
\leanok
\uses{Turing.Cfg, Turing.RecordAt, Turing.TableAt, Turing.TableWF, Turing.fieldOffset, Turing.joinFields, Turing.rt_fields_slice, Turing.tableFields, Turing.tableOffset}
\difficulty{4}
Let a configuration of a machine with $K$ work tapes have tape $r$ storing a table with
header field $h$ and entry fields $e_0,e_1,\ldots$, and suppose this table is
represented on the tape with separators between fields. Assume $h$ is nonempty, every
entry has length $m$, and every field within each entry is nonempty. If the head of tape
$r$ is at offset $o$ plus the starting offset of entry $e_j$, where $j$ is an entry
index and $j$ is less than the number of entries, then for every positive integer
$m_0\le m$, the tape stores the first $m_0$ fields of $e_j$ as a record beginning at its
head, with the record's prescribed separator layout.
\end{theorem}

\begin{definition}[Table rewind machine]
\lean{Turing.tableRewindTM}
\label{Turing.tableRewindTM}
\leanok
\uses{Turing.MultiTapeTM, Turing.SeekPhase, Turing.seamCompTM, Turing.seekStartTM, Turing.skipBackTM}
For an alphabet $\Gamma$, a natural number $k$ of work tapes, and a work-tape index $i$,
define the table-rewind machine to be the seam composition of the backward-skip machine
and the start-seeking machine, using the designated done and run states, respectively.
\end{definition}

\begin{theorem}[Rewinding a table tape]
\lean{Turing.tableRewindTM_run_ofCfg}
\label{Turing.tableRewindTM_run_ofCfg}
\leanok
\uses{Turing.Cfg, Turing.Cfg.mapState, Turing.FinTM.bufferTape, Turing.MultiTapeTM.runFrom, Turing.MultiTapeTM.visitedByTapeHead, Turing.SeekPhase, Turing.fieldOffset, Turing.fieldOffset_length, Turing.fieldOffset_succ, Turing.fieldsTape, Turing.fieldsTape_field, Turing.rt_fields_slice, Turing.seamCompTM_firstReturn_ofCfg, Turing.seamCompTM_run_ofCfg, Turing.seamCompTM_visitedByTapeHead_ofCfg, Turing.seekStartTM, Turing.seekStartTM_run_ofCfg, Turing.skipBackTM, Turing.skipBackTM_run_ofCfg, Turing.tableRewindTM}
\difficulty{7}
Let $k$ be a number of work tapes, let $\Gamma$ be an alphabet, let $x$ be an input
word, and choose a tape $i$ among the $k$ work tapes. Let $ws$ be a nonempty list of
nonempty fields, and write $L$ for its total field offset. Suppose a configuration $d$
starts in the left component's run phase, and for every integer $p$ with $-2 \le p < L$,
tape $i$ at position $d_i-L+p$ contains the symbol prescribed by the fields tape at
position $p$. After $L+7$ steps of the table-rewind machine, the configuration is
exactly $d$ with state changed to the right component's done phase and tape $i$'s head
moved to $d_i-L$. Before that time the machine has not reached the right component's
done phase; throughout steps $t \le L+7$, the head on tape $i$ remains between $d_i-L-2$
and $d_i$, inclusive, and every other work-tape head remains at its initial position.
\end{theorem}

\begin{definition}[Recursive copy-machine states]
\lean{Turing.rt_copyState}
\label{Turing.rt_copyState}
\leanok
\uses{Turing.SeekPhase, Turing.SweepPhase}
For each natural number $m$, define a state type recursively: at $m=0$ it is a
singleton, at $m=1$ it consists of the three sweep phases, and at $m+2$ it is the
disjoint union of the sweep phases, two copies of the seek phases, the state type at
$m+1$, and two further copies of the seek phases.
\end{definition}

\begin{definition}[Initial copy-machine state]
\lean{Turing.rt_copyEntry}
\label{Turing.rt_copyEntry}
\leanok
\uses{Turing.rt_copyState}
For each natural number $m$, the initial state is the unique state when $m=0$, the
forward sweep phase when $m=1$, and the forward sweep phase in the first summand of the
state type when $m\ge 2$.
\end{definition}

\begin{definition}[Final copy-machine state]
\lean{Turing.rt_copyExit}
\label{Turing.rt_copyExit}
\leanok
\uses{Turing.rt_copyState}
For each natural number $m$, the final state is the unique state when $m=0$, the exit
phase when $m=1$, and the exit phase in the last seek-phase summand of the state type
when $m\ge 2$.
\end{definition}

\begin{definition}[Recursive copy machine]
\lean{Turing.rt_copyMachine}
\label{Turing.rt_copyMachine}
\leanok
\uses{Turing.MultiTapeTM, Turing.copyTM, Turing.rt_copyEntry, Turing.rt_copyExit, Turing.rt_copyState, Turing.seamCompTM, Turing.skipBackTM, Turing.skipTM}
For each natural number $m$ and symbol alphabet $\Gamma$, define a deterministic Turing
machine with two work tapes and the recursively specified state type. At $m=0$ it stays
forever in its unique state without changing either work tape; at $m=1$ it copies from
the first work tape to the second; for $m\ge 2$ it composes a copy, two tape-skipping
machines, a recursive copy machine, and two returning machines, switching between them
at their designated exit states.
\end{definition}

\begin{theorem}[Base cases for copying joined fields]
\lean{Turing.rt_copy_run_base}
\label{Turing.rt_copy_run_base}
\leanok
\uses{Turing.Cfg, Turing.FinTM.bufferTape, Turing.MultiTapeTM.runFrom, Turing.RecordAt, Turing.copyTM, Turing.copyTM_run_ofCfg, Turing.fieldOffset, Turing.fieldsTape, Turing.fieldsTape_field, Turing.joinFields, Turing.rt_copyEntry, Turing.rt_copyExit, Turing.rt_copyMachine, Turing.rt_copyState}
\difficulty{5}
For every type $\Gamma$, natural number $m\le 1$, there is a natural-number constant $C$
such that for every input word $x$, list of words $f_1,\ldots,f_m$ over $\Gamma$, and
starting configuration $d$ of the specified two-tape machine, the following holds. If
the machine starts in its designated entry state, tape 0 stores the joined fields
throughout offsets $-1$ through their length $L$, and tape 1 is blank throughout those
offsets, then it reaches its exit state at some time $t\le C(L+1)$, avoiding that state
at every earlier time; both heads stay within distance $C(L+1)$ of their initial
positions up to and including time $t$, and at time $t$ tape 1 stores the joined fields
in the interval of length $L$ beginning at its initial head position, with all other
tape contents and configuration data unchanged.
\end{theorem}

\begin{theorem}[Sequential composition across a state seam]
\lean{Turing.rt_copy_seam}
\label{Turing.rt_copy_seam}
\leanok
\uses{Turing.Cfg, Turing.Cfg.mapState, Turing.MultiTapeTM, Turing.MultiTapeTM.runFrom, Turing.seamCompTM, Turing.seamCompTM_firstReturn_ofCfg, Turing.seamCompTM_run_ofCfg, Turing.seamCompTM_visitedByTapeHead_ofCfg}
\difficulty{4}
Let $\Gamma,S_1,S_2$ be types, let $k$ be a natural number, and let $x$ be an input word
over $\Gamma$. Let $M_1,M_2$ be machines with $k$ work tapes and state sets $S_1,S_2$,
respectively, with decidable equality on $S_1$. Suppose $M_1$ runs from configuration
$a$ to configuration $b$ in $t_1$ steps, with state $q_1$ at $b$ and no earlier
occurrence of that state; suppose every work-tape head during this run stays within
distance $B$ of a specified origin. Suppose also that $M_2$, started from $b$ with its
state renamed to the entry state, reaches configuration $c$ in $t_2$ steps, with state
$q_2$ at $c$ and no earlier occurrence of that state, and that its work-tape heads stay
within distance $B$ of the same origin. Then the machine that switches from $M_1$ to
$M_2$ at state $q_1$ reaches the state-renamed configuration $c$ after $t_1+1+t_2$
steps, does not reach state $q_2$ in the second component earlier, and keeps every
work-tape head within distance $B$ of the origin through that time.
\end{theorem}

\begin{theorem}[Copying joined fields]
\lean{Turing.rt_copy_run}
\label{Turing.rt_copy_run}
\leanok
\uses{Turing.Cfg, Turing.Cfg.mapState, Turing.FinTM.bufferTape, Turing.MultiTapeTM, Turing.MultiTapeTM.runFrom, Turing.RecordAt, Turing.SeekPhase, Turing.SweepPhase, Turing.copyTM, Turing.copyTM_run_ofCfg, Turing.fieldOffset, Turing.fieldsTape, Turing.fieldsTape_field, Turing.joinFields, Turing.joinFields_cons_cons, Turing.rt_copyEntry, Turing.rt_copyExit, Turing.rt_copyMachine, Turing.rt_copyState, Turing.rt_copy_run_base, Turing.rt_copy_seam, Turing.rt_fields_slice, Turing.seamCompTM, Turing.skipBackTM, Turing.skipBackTM_run_ofCfg, Turing.skipTM, Turing.skipTM_run_ofCfg}
\difficulty{7}
For every type $\Gamma$ and natural number $m$, there is a natural-number constant $C$
such that the following holds for every input word $x$, list of words $f_1,\ldots,f_m$
over $\Gamma$, and starting configuration $d$ of the specified two-tape machine. If tape
0 stores the joined fields throughout offsets $-1$ through the joined-fields length,
tape 1 is blank throughout that interval, and the machine starts in its designated entry
state, then it reaches its exit state at some time $t\le C(L+1)$, where $L$ is the
joined-fields length. The exit state is avoided at every earlier time, both work-tape
heads remain within distance $C(L+1)$ of their initial positions throughout the run, and
at the final time tape 1 stores the joined fields in the interval of length $L$
beginning at its initial head position while all other tape contents and configuration
data are unchanged.
\end{theorem}

\begin{theorem}[Machine copying joined fields]
\lean{Turing.exists_fieldsCopyTM}
\label{Turing.exists_fieldsCopyTM}
\leanok
\uses{Turing.Cfg, Turing.MultiTapeTM, Turing.MultiTapeTM.runFrom, Turing.RecordAt, Turing.fieldsTape, Turing.joinFields, Turing.rt_copyEntry, Turing.rt_copyExit, Turing.rt_copyMachine, Turing.rt_copyState, Turing.rt_copy_run}
\difficulty{1}
For every alphabet $\Gamma$ and natural number $m$, there exist a finite state set, a
deterministic Turing machine over $\Gamma$ with two work tapes, entry and exit states,
and a constant $C$. For every input word, every list $f_s$ of $m$ fields, and every
configuration $d$ in the entry state whose tape $0$ records the joined fields and whose
tape $1$ is blank on the cells from one position before the record through its last
position, the machine reaches its exit state at some time $t \le C(L+1)$, where $L$ is
the length of the joined record. It does not visit the exit state earlier, every
work-tape head stays within distance $C(L+1)$ of its initial position through time $t$,
and at time $t$ tape $1$ contains the joined fields on the interval from its initial
head position up to but excluding that position plus $L$; all other tape cells and the
remaining configuration data are unchanged.
\end{theorem}

\begin{definition}[Find-machine phases]
\lean{Turing.rt_findPhase}
\label{Turing.rt_findPhase}
\leanok
Given a natural number $m$ and a type $A$, define a phase type with a read phase, scan
phases indexed by $j\in\{0,\ldots,m-1\}$, skip phases indexed by such a $j$ and an
element of $A$, and an exit phase.
\end{definition}

\begin{definition}[Control states for find phases]
\lean{Turing.rt_findControl}
\label{Turing.rt_findControl}
\leanok
\uses{Turing.ScanPhase, Turing.SeekPhase, Turing.rt_findPhase}
Given a natural number $m$ and a type $A$, assign control states to each find-machine
phase: the read phase has control states in
$\mathrm{Option}(\{\mathrm{false},\mathrm{true}\})$, each scan phase has a scan phase
state, each skip phase has a seek phase state, and the exit phase has a singleton
control state.
\end{definition}

\begin{definition}[Find-machine component machines]
\lean{Turing.rt_findParts}
\label{Turing.rt_findParts}
\leanok
\uses{Turing.MultiTapeTM, Turing.inputMoveAction, Turing.rt_findControl, Turing.rt_findPhase, Turing.scanFold, Turing.scanKeep, Turing.scanTM, Turing.skipTM}
Let $\Gamma$ and $A$ be types, let $m$ be a natural number, let
$\delta:A\to\mathrm{Option}(\Gamma)\to A$, and choose an initial element $a_0\in A$. For
each find-machine phase, define a two-work-tape Turing machine over $\Gamma$ with that
phase's control-state type: the read phase records whether the first work-tape cell is
occupied, each scan phase performs a forward scan using the state update
$a\mapsto\delta(a,\mathrm{some}(\gamma))$ while preserving scanned symbols, each skip
phase advances along the first tape, and the exit phase remains at its singleton control
state.
\end{definition}

\begin{definition}[Find-machine phase redirection]
\lean{Turing.rt_findRedirect}
\label{Turing.rt_findRedirect}
\leanok
\uses{Turing.rt_findControl, Turing.rt_findPhase}
Let $\Gamma$ and $A$ be types, let $m>0$, let $\delta:A\to\mathrm{Option}(\Gamma)\to A$,
choose $a_0\in A$, and let $\mathrm{acc}:A\to\{\mathrm{false},\mathrm{true}\}$. Define a
partial redirection from a find-machine phase and its control state to a new phase and
control state as follows: an empty read control halts, a false read control goes to
exit, and a true read control begins a forward scan at index $0$ with state $a_0$; a
completed scan either skips to the next index with updated state
$\delta(a,\mathrm{none})$, goes to exit if it is the last index and $\mathrm{acc}(a)$
holds, or skips at the last index with state $a$; a completed skip advances to the next
scan index when one remains and otherwise returns to the read phase. All other
phase-control combinations halt.
\end{definition}

\begin{definition}[Round-finding machine construction]
\lean{Turing.rt_findRound}
\label{Turing.rt_findRound}
\leanok
\uses{Turing.MultiTapeTM, Turing.phaseHostTM, Turing.rt_findControl, Turing.rt_findParts, Turing.rt_findPhase, Turing.rt_findRedirect}
For a positive integer $m$, types $A$ and $\Gamma$, a transition function
$\delta:A\to(\mathrm{Option}\,\Gamma)\to A$, an initial element $a_0\in A$, and a
Boolean-valued function $\mathrm{acc}:A\to\{\mathrm{false},\mathrm{true}\}$, the
round-finding construction defines a deterministic Turing machine with two work tapes,
alphabet $\Gamma$, and states drawn from a phase together with phase-dependent control
data. Its initial state is the read phase with control value none.
\end{definition}

\begin{theorem}[Two-way scan of a word]
\lean{Turing.rt_find_scan}
\label{Turing.rt_find_scan}
\leanok
\uses{Turing.Cfg, Turing.FinTM.bufferTape, Turing.MultiTapeTM.runFrom, Turing.ScanPhase, Turing.scanFold, Turing.scanKeep, Turing.scanResult_fold, Turing.scanTM, Turing.scanTM_run_ofCfg}
\difficulty{4}
Let $\delta:A\times\mathrm{Option}(\Gamma)\to A$, let $a\in A$, and let $w$ be a word
over $\Gamma$. Suppose a two-work-tape configuration $d$ has forward scan state $a$, and
tape 0, at each offset $p$ from $-1$ through $|w|$ relative to its head, contains the
corresponding symbol of $w$ when $0\le p<|w|$ and is blank otherwise. After $2|w|+2$
steps of the scanning machine, its state is the completed state whose accumulator is the
result of folding $\delta$ over $w$ from $a$, with all other configuration data
unchanged; at every earlier time its state is not a completed state. At every time up to
and including $2|w|+2$, the head of tape 0 lies between its initial position minus one
and its initial position plus $|w|$, inclusive, and every other work-tape head remains
at its initial position.
\end{theorem}

\begin{theorem}[Fold across joined fields]
\lean{Turing.rt_find_fold_cons}
\label{Turing.rt_find_fold_cons}
\leanok
\uses{Turing.joinFields, Turing.joinFields_cons_cons}
Let $\Gamma$ and $A$ be types, let $\delta:A\to\mathrm{Option}(\Gamma)\to A$, and
let $a\in A$, $u,v$ be lists of elements of $\Gamma$, and $fs$ be a list of such lists.
Folding the joined fields $u,v,fs$ from $a$ gives the same result as folding the joined
fields $v,fs$ from $\delta(b,\mathrm{none})$, where $b$ is obtained by folding $u$
from $a$ and supplying each element of $u$ to $\delta$ as a present symbol.
\end{theorem}

\begin{theorem}[Initial read phase]
\lean{Turing.rt_find_read}
\label{Turing.rt_find_read}
\leanok
\uses{Turing.Cfg, Turing.Cfg.mapState, Turing.MultiTapeTM.runFrom, Turing.MultiTapeTM.runFrom_succ_eq_step', Turing.MultiTapeTM.runFrom_zero, Turing.moveInputPos_zero, Turing.phaseHostTM_run, Turing.rt_findParts, Turing.rt_findPhase, Turing.rt_findRedirect, Turing.rt_findRound, Turing.step_of_inputMoveAction}
\difficulty{5}
Let $m>0$, let $\delta:A\times(\Gamma\cup\{\text{blank}	o A$ be the state update
function, and let $a_0\in A$ and $\mathrm{acc}:A\to\{\mathrm{false},\mathrm{true}\}$.
For a two-work-tape configuration $d$ on input $x$ whose control state is the initial
read control, after one step of the constructed machine its control is the read phase
with a bit recording whether tape 0 has a symbol under its head. After two steps its
control is the initial forward scan state on tape 0 if that cell is occupied, and is the
exit state if it is blank; in both equalities all other configuration data are as in
$d$.
\end{theorem}

\begin{theorem}[Host behavior during a scan phase]
\lean{Turing.rt_find_scan_phase}
\label{Turing.rt_find_scan_phase}
\leanok
\uses{Turing.Cfg, Turing.Cfg.mapState, Turing.FinTM.bufferTape, Turing.MultiTapeTM.runFrom, Turing.ScanPhase, Turing.phaseHostTM_run, Turing.rt_findControl, Turing.rt_findParts, Turing.rt_findPhase, Turing.rt_findRedirect, Turing.rt_findRound, Turing.rt_find_scan, Turing.scanFold, Turing.scanKeep, Turing.scanTM}
\difficulty{4}
Let $m>0$, let $\delta:A\times\mathrm{Option}(\Gamma)\to A$, choose $a_0\in A$ and an
acceptance predicate $\mathrm{acc}:A\to\{\mathrm{false},\mathrm{true}\}$, and let $j$ be
an index with $0\le j<m$. Let $a\in A$, let $w$ be a word over $\Gamma$, and suppose a
two-work-tape configuration $d$ has forward scan state $a$ and tape 0 contains the
encoding of $w$ at all offsets from $-1$ through $|w|$. Let $q$ be a combined phase and
control state such that the host's redirect from scan phase $j$ with completed
accumulator equal to the result of folding $\delta$ over $w$ from $a$ leads to $q$. For
every $u\le 2|w|+2$, the host machine started with $d$ in phase $j$ agrees with the
standalone scan machine for $u$ steps, with its control state tagged by phase $j$. One
further step takes the host configuration to $d$ with its state replaced by $q$.
\end{theorem}

\begin{theorem}[Host behavior during a skip phase]
\lean{Turing.rt_find_skip_phase}
\label{Turing.rt_find_skip_phase}
\leanok
\uses{Turing.Cfg, Turing.Cfg.mapState, Turing.FinTM.bufferTape, Turing.MultiTapeTM.runFrom, Turing.SeekPhase, Turing.phaseHostTM_run, Turing.rt_findControl, Turing.rt_findParts, Turing.rt_findPhase, Turing.rt_findRedirect, Turing.rt_findRound, Turing.skipTM, Turing.skipTM_run_ofCfg}
\difficulty{4}
Let $m>0$, let $\delta:A\times\mathrm{Option}(\Gamma)\to A$, choose $a_0\in A$ and an
acceptance predicate $\mathrm{acc}:A\to\{\mathrm{false},\mathrm{true}\}$, and let $j<m$
and $a\in A$. Suppose $w$ is a word over $\Gamma$ and a two-work-tape configuration $d$
is in the skip machine's run state, with tape 0 containing the encoding of $w$ at every
offset from $0$ through $|w|$. Let $q$ be a combined phase and control state such that
the host's redirect from skip phase $j$ in its completed state leads to $q$. For every
$u\le |w|+1$, the host started with $d$ in skip phase $j$ agrees with the standalone
skip machine for $u$ steps, with its control state tagged by phase $j$. After one
further step, its state is $q$ and its tape 0 head is at its initial position plus
$|w|+1$; all other configuration data agree with $d$.
\end{theorem}

\begin{theorem}[Accepting final scan]
\lean{Turing.rt_find_last_accept}
\label{Turing.rt_find_last_accept}
\leanok
\uses{Turing.Cfg, Turing.Cfg.mapState, Turing.FinTM.bufferTape, Turing.MultiTapeTM.runFrom, Turing.MultiTapeTM.runFrom_state_ne_none_of_le, Turing.ScanPhase, Turing.rt_findPhase, Turing.rt_findRedirect, Turing.rt_findRound, Turing.rt_find_scan, Turing.rt_find_scan_phase, Turing.scanFold, Turing.scanKeep, Turing.scanTM}
\difficulty{5}
Let $m>0$, let $\Gamma$ and $A$ be types, let $\delta:A\to\mathrm{Option}(\Gamma)\to A$,
and fix $a_0\in A$ and $acc:A\to\{\mathrm{false},\mathrm{true}\}$. Let $j\in\mathrm{Fin}(m)$
be the final field index, so $j+1=m$. Let $a\in A$ and let $w$ be a word over $\Gamma$,
with $acc$ true on the state obtained by folding $w$ from $a$. Let $d$ be a configuration
on input $x$ in the forward scan state with value $a$, whose tape $0$ agrees with the
buffer for $w$ at every relative integer position from $-1$ through $|w|$. After
$2|w|+3$ steps, the round machine from the scan phase for $j$ reaches $d$ mapped to the
exit state, with all other configuration data unchanged. At every earlier time $u<2|w|+3$,
its state exists and is neither the read state with no stored Boolean value nor the exit
state. Throughout every time $u\le2|w|+3$, the head on tape $0$ stays between its initial
position minus one and its initial position plus $|w|$, inclusive; the head on tape $1$
remains fixed.
\end{theorem}

\begin{theorem}[Advance through a scanned field]
\lean{Turing.rt_find_advance}
\label{Turing.rt_find_advance}
\leanok
\uses{Turing.Cfg, Turing.Cfg.mapState, Turing.FinTM.bufferTape, Turing.MultiTapeTM.runFrom, Turing.MultiTapeTM.runFrom_add, Turing.ScanPhase, Turing.SeekPhase, Turing.rt_findControl, Turing.rt_findPhase, Turing.rt_findRedirect, Turing.rt_findRound, Turing.rt_find_scan_phase, Turing.rt_find_skip_phase, Turing.scanFold, Turing.scanKeep, Turing.scanTM, Turing.skipTM}
\difficulty{5}
Let $m>0$, let $\Gamma,A$ be types, and fix $\delta:A\to\mathrm{Option}(\Gamma)\to A$,
$a_0\in A$, and $acc:A\to\{\mathrm{false},\mathrm{true}\}$. Let $j\in\mathrm{Fin}(m)$,
let $a,b\in A$, let $w$ be a word over $\Gamma$, and let $d$ be a configuration in the
forward scan state with value $a$. Suppose tape $0$, relative to its head, contains the
buffer for $w$ at every integer position from $-1$ through $|w|$. Suppose the phase
redirect from scanning field $j$ after folding $w$ from $a$ leads to the skip phase
carrying $b$, and that the completed skip phase carrying $b$ redirects to a phase-control
pair $q$. After $3|w|+5$ steps, the round machine reaches the original configuration
with tape $0$'s head advanced by $|w|+1$ and its state mapped to $q$. At every earlier
time $u<3|w|+5$, its state is neither the read state with no stored Boolean value nor
the exit state.
\end{theorem}

\begin{theorem}[Run through the remaining fields]
\lean{Turing.rt_find_fields_run}
\label{Turing.rt_find_fields_run}
\leanok
\uses{Turing.Cfg, Turing.Cfg.mapState, Turing.FinTM.bufferTape, Turing.MultiTapeTM.runFrom, Turing.MultiTapeTM.runFrom_add, Turing.ScanPhase, Turing.fieldOffset, Turing.fieldOffset_add_length_le, Turing.fieldOffset_cons_succ, Turing.fieldsTape, Turing.fieldsTape_cons_shift, Turing.fieldsTape_field, Turing.joinFields, Turing.joinFields_cons_cons, Turing.rt_findPhase, Turing.rt_findRedirect, Turing.rt_findRound, Turing.rt_find_advance, Turing.rt_find_fold_cons, Turing.rt_find_last_accept}
\difficulty{6}
Let $m>0$, let $\Gamma,A$ be types, and fix $\delta:A\to\mathrm{Option}(\Gamma)\to A$,
$a_0\in A$, and $acc:A\to\{\mathrm{false},\mathrm{true}\}$. Let $fs$ be a list of
lists over $\Gamma$, and let $j\in\mathrm{Fin}(m)$ satisfy $j+|fs|=m$. Let $a\in A$ and
let $d$ be a configuration in the forward scan state with value $a$. Suppose tape $0$,
relative to its head, agrees with the tape formed by joining $fs$ at every integer
position from $-1$ through $|\mathrm{joinFields}(fs)|$. Then there is a time
$t\le5(|\mathrm{joinFields}(fs)|+1)$ such that at every time $u<t$ the run from $d$
mapped to scan phase $j$ is in neither the read state with no stored Boolean value nor
the exit state. At time $t$, it is the original configuration with tape $0$'s head
advanced by $\mathrm{fieldOffset}(fs,|fs|-1)$ if $acc$ accepts the fold of
$\mathrm{joinFields}(fs)$ from $a$, and by $\mathrm{fieldOffset}(fs,|fs|)$ otherwise;
its state is respectively the exit state or the read state with no stored Boolean value.
\end{theorem}

\begin{theorem}[Blank input transitions to exit]
\lean{Turing.rt_find_blank}
\label{Turing.rt_find_blank}
\leanok
\uses{Turing.Cfg, Turing.Cfg.mapState, Turing.MultiTapeTM.runFrom, Turing.rt_findPhase, Turing.rt_findRound, Turing.rt_find_read}
\difficulty{4}
Let $m>0$, let $\Gamma,A$ be types, and fix $\delta:A\to\mathrm{Option}(\Gamma)\to A$,
$a_0\in A$, and $acc:A\to\{\mathrm{false},\mathrm{true}\}$. Let $d$ be a configuration
on input $x$ in state $none$. If the cell under work-tape head $0$ is blank, then after
two steps the round machine, started from $d$ mapped to the read phase with no captured
value, reaches $d$ mapped to the exit state, with all other configuration data
unchanged. At every integer time $u$ with $0<u<2$, its state is neither the read state
with no stored Boolean value nor the exit state.
\end{theorem}

\begin{theorem}[Finding a Field in Joined Data]
\lean{Turing.rt_find_entry}
\label{Turing.rt_find_entry}
\leanok
\uses{Turing.Cfg, Turing.Cfg.mapState, Turing.FinTM.bufferTape, Turing.MultiTapeTM.runFrom, Turing.MultiTapeTM.runFrom_add, Turing.MultiTapeTM.runFrom_state_ne_none_of_le, Turing.ScanPhase, Turing.fieldOffset, Turing.fieldsTape, Turing.fieldsTape_field, Turing.joinFields, Turing.rt_findControl, Turing.rt_findPhase, Turing.rt_findRound, Turing.rt_find_fields_run, Turing.rt_find_read}
\difficulty{6}
Let $m>0$, let $\Gamma$ and $A$ be types, let $\delta:A\to\mathrm{Option}(\Gamma)\to A$,
and choose $a_0\in A$ and $acc:A\to\{\mathrm{false},\mathrm{true}\}$. Let $fs$ be a list
of exactly $m$ nonempty finite lists over $\Gamma$, and let $d$ be a two-work-tape
configuration on input $x$ in state $none$. Suppose the first work tape of $d$, at every
offset from $-1$ through the length of the joined encoding of $fs$, agrees with that
encoding. Then the find machine has a positive time $t\le 7(|joinFields(fs)|+1)$ at
which it reaches the otherwise unchanged configuration except that the first work-tape
head is advanced by $fieldOffset(fs,j)$, where $j=m-1$ if $acc$ accepts the fold of
$\delta$ over the joined fields starting at $a_0$, and $j=m$ otherwise. At every time
$u$ with $0<u<t$, its state is neither the initial read state nor the exit state; at
time $t$ it is in the exit state if the fold is accepted, and otherwise in the initial
read state.
\end{theorem}

\begin{definition}[Round-finding clock]
\lean{Turing.rt_findClock}
\label{Turing.rt_findClock}
\leanok
\uses{Turing.Cfg, Turing.MultiTapeTM.runFrom, Turing.rt_findControl, Turing.rt_findPhase, Turing.rt_findRound}
Let $m>0$, let $A$ and $\Gamma$ be types, let $\delta:A\to\mathrm{Option}(\Gamma)\to A$,
let $a_0\in A$, and let $\mathrm{acc}:A\to\{\mathrm{false},\mathrm{true}\}$. For an
input list $x$ over $\Gamma$ and a configuration $d$ of the round-finding machine,
define its clock to be the least positive time $t$ at which the machine's state is
either the read phase with control value none or the exit phase with its unit control
value; if no such time exists, define the clock to be $1$.
\end{definition}

\begin{theorem}[Least event time]
\lean{Turing.rt_findClock_eq}
\label{Turing.rt_findClock_eq}
\leanok
\uses{Turing.Cfg, Turing.MultiTapeTM.runFrom, Turing.rt_findClock, Turing.rt_findControl, Turing.rt_findPhase, Turing.rt_findRound}
\difficulty{4}
Let $m>0$, let $\Gamma,A$ be types, and fix $\delta:A\to\mathrm{Option}(\Gamma)\to A$,
$a_0\in A$, and $acc:A\to\{\mathrm{false},\mathrm{true}\}$. Let $d$ be a configuration
of the round machine whose state records a phase and its control value. Suppose $t>0$ and
after $t$ steps the machine is in either the read state with no stored Boolean value or
the exit state. Suppose that at every positive time $u<t$ it is in neither state. Then
the proof clock of $d$, defined as the first positive time it reaches either state, is $t$.
\end{theorem}

\begin{theorem}[Return to the round boundary]
\lean{Turing.rt_find_round_entry}
\label{Turing.rt_find_round_entry}
\leanok
\uses{Turing.Cfg, Turing.Cfg.mapState, Turing.MultiTapeTM.runFrom, Turing.fieldOffset, Turing.fieldsTape, Turing.joinFields, Turing.rt_findClock, Turing.rt_findClock_eq, Turing.rt_findControl, Turing.rt_findPhase, Turing.rt_findRound, Turing.rt_find_entry}
\difficulty{5}
Let $m>0$, let $\delta:A\times\mathrm{Option}(\Gamma)\to A$, choose $a_0\in A$ and an
acceptance predicate $\mathrm{acc}:A\to\{\mathrm{false},\mathrm{true}\}$. Let
$f_1,\ldots,f_m$ be nonempty words over $\Gamma$, and form $J$ by joining them in order
with one blank symbol between consecutive words. Suppose a two-work-tape configuration
$c$ is at the round's read control and, on tape 0, its cells at offsets $p$ from the
head agree with the tape encoding of $J$ for every integer $p$ from $-1$ through $|J|$.
Let $t$ be the first positive time at which the machine's control is either the read
control or the exit control. Then $0<t\le 7(|J|+1)$; at every intermediate time $0<u<t$
its control is neither of these controls. At time $t$, the configuration agrees with $c$
except that its control is the exit control if $\mathrm{acc}$ accepts the state obtained
by folding $\delta$ over $J$ from $a_0$, and otherwise is the read control; the head of
tape 0 has moved by the sum of the lengths of the first $m-1$ words plus their
separators in the accepted case, and by $|J|+1$ in the other case.
\end{theorem}

\begin{theorem}[Linear-time table row search]
\lean{Turing.exists_tableFindTM}
\label{Turing.exists_tableFindTM}
\leanok
\uses{Turing.Cfg, Turing.FinTM.bufferTape, Turing.MultiTapeTM, Turing.MultiTapeTM.runFrom, Turing.TableAt, Turing.TableWF, Turing.joinFields, Turing.tableFields, Turing.tableOffset}
Let $\Gamma$ be a type with decidable equality, let $e:\{0,1\}\hookrightarrow\Gamma$ be
an injection, and let $m$ be a positive integer. For every finite type $A$ with
decidable equality, transition map $\delta:A\times(\Gamma\cup\{\text{blank}\	o A$,
initial state $a_0\in A$, and predicate $\mathrm{acc}:A\to\{0,1\}$, there are a natural
number $K$, a finite state type $P$, a deterministic machine $R$ with $K+2$ work tapes
and distinct states $\mathrm{yes},\mathrm{no}\in P$, an entry state, and a constant
$C\in\mathbb N$ with the following property. For any input word $y$, integer offset $o$,
header $h$, list of rows $es$, and Boolean word $w$, suppose $h$ is nonempty, every row
has exactly $m$ nonempty fields, and the number of rows is strictly less than the
integer whose binary digits are the entries of $w$. If the machine starts in the entry
state with tape 0's head at $o$, tape 0 contains the joined header and rows at offset
$o$, tape 1 contains the encoding of $w$ starting at its head, and every tape numbered
at least 2 is blank, then there is a time $t\le C(L+|w|+1)$, where $L$ is the length of
the joined table, at which the machine is in the yes state exactly when some row, in
list order, yields an accepting state after its joined fields are processed from $a_0$
by repeatedly applying $\delta$ to the current state and next symbol; if no row does, it
is in the no state. At every earlier time it is in neither the yes nor the no state, and
through time $t$ each work-tape head remains within distance $C(L+|w|+1)$ of its
starting position. At time $t$, in the yes case tape 0's head is at $o$ plus the offset
of the first accepting row; in the no case the final configuration equals the initial
configuration except that its state is the no state.
\end{theorem}

\begin{theorem}[Bounded reverse field copying]
\lean{Turing.rt_fieldsCopy_reverse}
\label{Turing.rt_fieldsCopy_reverse}
\leanok
\uses{Turing.Cfg, Turing.MultiTapeTM, Turing.MultiTapeTM.runFrom, Turing.RecordAt, Turing.embedEmitCfg, Turing.embedEmitCfg_selected_pos, Turing.embedEmitCfg_selected_tape, Turing.embedEmitTM, Turing.embedEmitTM_runFrom, Turing.exists_fieldsCopyTM, Turing.fieldsTape, Turing.joinFields}
\difficulty{5}
For every natural number $m$, there are a finite state set $P$ with decidable equality,
a two-work-tape machine $R$ over an arbitrary symbol type $\Gamma$, designated states
$entry,exit\in P$, and a natural number $C$ such that the following holds. For every
input word $x$, list of fields $fs$ of length $m$, and configuration $d$ on input $x$,
suppose $d$ begins in state $entry$, tape $1$ contains the joined fields in the
specified field-tape representation at offsets from $-1$ through the total joined length
relative to its head, and tape $0$ is blank throughout those offsets. Then in some time
$t\le C(L+1)$, where $L$ is the length of the joined fields, the machine reaches a
configuration equal to $d$ except that its state is $exit$ and tape $0$ contains the
joined fields in field-tape representation on the interval of length $L$ starting at its
initial head position. It does not visit $exit$ earlier, and both tape heads remain
within distance $C(L+1)$ of their initial positions through time $t$.
\end{theorem}

\begin{theorem}[Copying a record to the table]
\lean{Turing.rt_append_copy}
\label{Turing.rt_append_copy}
\leanok
\uses{Turing.Cfg, Turing.Cfg.mapState, Turing.MultiTapeTM, Turing.MultiTapeTM.runFrom, Turing.RecordAt, Turing.SeekPhase, Turing.TableAt, Turing.fieldOffset, Turing.fieldOffset_length, Turing.fieldsTape, Turing.joinFields, Turing.rt_copy_seam, Turing.rt_fieldsCopy_reverse, Turing.seamCompTM, Turing.seekEndTM, Turing.seekEndTM_run_ofCfg, Turing.tableFields, Turing.tableOffset, Turing.tableOffset_length}
\difficulty{6}
For every $m\in\mathbb{N}$, there are a finite state set $P$, a two-work-tape machine
$R$, entry and exit states, and a constant $C\in\mathbb{N}$ such that the following
holds. Given a header word $h$, a list $es$ of lists of words, and a list $e$ of words
of length $m$, assume every field in the table formed from $h$ and the flattened $es$ is
nonempty. If a configuration starts in the entry state, tape 0 has its head at $o$ and
stores the table encoding with blank cells from its end onward, and tape 1 stores the
joined encoding of $e$ at its head, then there is a time $t\le C(L+E+1)$, where $L$ and
$E$ are the lengths of the table and record encodings, at which the machine is in the
exit state and tape 0 has appended the encoding of $e$ at the offset following all table
fields. No earlier time below $t$ reaches the exit state, and up to time $t$ both heads
remain within distance $C(L+E+1)$ of their initial positions; the final tape contents
elsewhere, as well as the other configuration data, are unchanged.
\end{theorem}

\begin{theorem}[Appending a row to a table]
\lean{Turing.exists_tableAppendTM}
\label{Turing.exists_tableAppendTM}
\leanok
\uses{Turing.Cfg, Turing.MultiTapeTM, Turing.MultiTapeTM.runFrom, Turing.RecordAt, Turing.TableAt, Turing.fieldsTape, Turing.joinFields, Turing.tableFields}
For every natural number $m$ and fixed suffix of fields, there exist a finite-state
two-work-tape machine $R$, entry and exit states, and a constant $C$. Suppose a table
consists of a nonempty header $h$ and rows $es$, every field in the table is nonempty,
the appended row $e$ has exactly $m$ fields, tape 0 at offset $o$ contains the table
encoding with blank cells beyond its end, and tape 1 contains the encoding of $e$ at its
head. Then $R$ reaches, in at most $C(L+E+1)$ steps, the configuration obtained by
appending $e$ followed by the fixed suffix to the table, where $L$ is the encoded length
of the original table and $E$ is the encoded length of $e$. It does not enter the exit
state earlier, and throughout the run each work-tape head moves by at most $C(L+E+1)$
from its initial position.
\end{theorem}

\begin{definition}[Fixed-width Boolean successor word]
\lean{Turing.rt_succWord}
\label{Turing.rt_succWord}
\leanok
\uses{Turing.incFixed}
For a finite list $w$ of Boolean digits, return its fixed-width increment when that
increment is defined; otherwise return the list of $|w|$ false digits followed by a true
digit.
\end{definition}

\begin{theorem}[Successor word on a leading one]
\lean{Turing.rt_succWord_cons}
\label{Turing.rt_succWord_cons}
\leanok
\uses{Turing.incFixed, Turing.rt_succWord}
\difficulty{3}
For every Boolean word $w$, the successor word of the word formed by placing a one
before $w$ is a zero followed by the successor word of $w$.
\end{theorem}

\begin{theorem}[Binary successor of a natural number]
\lean{Turing.rt_succWord_bits}
\label{Turing.rt_succWord_bits}
\leanok
\uses{Turing.incFixed, Turing.rt_succWord, Turing.rt_succWord_cons}
\difficulty{4}
For every natural number $n$, the successor word associated with the binary expansion of
$n$ is the binary expansion of $n+1$.
\end{theorem}

\begin{theorem}[Length preservation under fixed-width increment]
\lean{Turing.rt_incFixed_length}
\label{Turing.rt_incFixed_length}
\leanok
\uses{Turing.incFixed}
\difficulty{4}
For Boolean words $w$ and $v$, if the fixed-width binary increment of $w$ returns $v$,
then $v$ and $w$ have the same length.
\end{theorem}

\begin{definition}[Boundary-writing machine]
\lean{Turing.rt_boundaryTM}
\label{Turing.rt_boundaryTM}
\leanok
\uses{Turing.MultiTapeTM, Turing.inputMoveAction, Turing.tapeAction}
Let $k$ be a positive integer, choose one of its work tapes, and fix a symbol $b$ in an
alphabet $\Gamma$. The resulting $k$-work-tape machine has four states, starts in state
$0$, moves the chosen work-tape head left in state $0$, writes $b$ in the scanned cell
and moves right in state $1$, then moves right in state $2$; in every other state it
makes no work-tape change and remains in state $3$. Its input head stays put throughout,
and it emits no output.
\end{definition}

\begin{theorem}[Three-step boundary operation]
\lean{Turing.rt_boundary_run}
\label{Turing.rt_boundary_run}
\leanok
\uses{Turing.Cfg, Turing.MultiTapeTM.runFrom, Turing.MultiTapeTM.step, Turing.rt_boundaryTM, Turing.step_of_tapeAction_move, Turing.step_of_tapeAction_write}
\difficulty{5}
Let $k\in\mathbb{N}$, let $\Gamma$ be a symbol alphabet, let $x$ be an input word, and
let $i$ be one of $k$ work tapes. For any symbol $b\in\Gamma$ and configuration $d$
whose state is $0$, the boundary machine reaches after three steps the configuration
with state $3$, with $b$ written on tape $i$ at the cell one left of its initial head
position and that head moved one cell to the right. State $3$ is not reached at any
earlier time, and through time three the selected head stays between one cell left and
one cell right of its initial position while every other work-tape head remains fixed.
\end{theorem}

\begin{definition}[Append-symbol machine]
\lean{Turing.rt_appendSymbolTM}
\label{Turing.rt_appendSymbolTM}
\leanok
\uses{Turing.MultiTapeTM, Turing.SeekPhase, Turing.rt_boundaryTM, Turing.seamCompTM, Turing.skipBackTM, Turing.skipTM}
Let $k$ be a positive integer, choose one of its work tapes, and fix a symbol $b$ in an
alphabet $\Gamma$. The resulting machine is the sequential composition of a machine that
skips forward on the selected tape, a boundary-writing machine that writes $b$, and a
machine that skips back; the first and second stages are joined at their specified exit
states, and the composite starts in the first stage's initial state.
\end{definition}

\begin{theorem}[Bounded construction of an appended record]
\lean{Turing.rt_appendSymbol_run}
\label{Turing.rt_appendSymbol_run}
\leanok
\uses{Turing.Cfg, Turing.Cfg.mapState, Turing.FinTM.bufferTape, Turing.FinTM.bufferTape_append, Turing.MultiTapeTM.runFrom, Turing.SeekPhase, Turing.rt_appendSymbolTM, Turing.rt_boundaryTM, Turing.rt_boundary_run, Turing.rt_copy_seam, Turing.seamCompTM, Turing.skipBackTM, Turing.skipBackTM_run_ofCfg, Turing.skipTM, Turing.skipTM_run_ofCfg}
\difficulty{6}
For every nonnegative integer $m$, there are a finite state set $P$, distinctness
decidable on $P$, a two-work-tape machine $R$, entry and exit states, and a constant
$C\in\mathbb{N}$ with the following property. Let $h$ be a header word, let $es$ be a
list of lists of words, and let $e$ be a list of words of length $m$, with every word in
the table formed by $h$ followed by the flattened entries of $es$ nonempty. If a
configuration $d$ starts in the entry state, tape 0 has head position $o$ and contains
that table in its joined-field encoding, tape 0 is blank at every position at or beyond
the table's end, and tape 1 contains the joined-field encoding of $e$ at its head, then
for some time $t\le C(L+E+1)$, where $L$ and $E$ are the lengths of those two joined
encodings, running $R$ for $t$ steps reaches the exit state and leaves every work-tape
head within distance $C(L+E+1)$ of its starting position. The exit is not reached
earlier; at time $t$, tape 0 contains the encoding of $e$ beginning at the table offset
after all its fields, with the tape head at that offset, while every other tape cell
agrees with its initial content.
\end{theorem}

\begin{definition}[Successor-machine phases]
\lean{Turing.rt_succPhase}
\label{Turing.rt_succPhase}
\leanok
Introduce a three-element phase type with phases for incrementing, extending, and
exiting.
\end{definition}

\begin{definition}[Successor-phase control data]
\lean{Turing.rt_succControl}
\label{Turing.rt_succControl}
\leanok
\uses{Turing.FlagPhase, Turing.SeekPhase, Turing.rt_succPhase}
For the incrementing phase, the control data is a flag phase; for the extending phase,
it is a disjoint union of a seek phase and a disjoint union of a four-state control and
another seek phase; for the exiting phase, it consists of the single unit value.
\end{definition}

\begin{definition}[Successor-machine components]
\lean{Turing.rt_succParts}
\label{Turing.rt_succParts}
\leanok
\uses{Turing.MultiTapeTM, Turing.incrementTM, Turing.inputMoveAction, Turing.rt_appendSymbolTM, Turing.rt_succControl, Turing.rt_succPhase}
Let $k$ be a positive integer and choose one of its work tapes. For each phase, assign a
machine with the corresponding phase-dependent control data: the incrementing phase uses
the Boolean increment machine, the extending phase uses the append-symbol machine with
symbol true, and the exiting phase uses a one-state machine that leaves the input head
and all work tapes unchanged, emits no output, and stays in its sole state.
\end{definition}

\begin{definition}[Phase-dependent redirect map]
\lean{Turing.rt_succRedirect}
\label{Turing.rt_succRedirect}
\leanok
\uses{Turing.rt_succControl, Turing.rt_succPhase}
Define a partial map between phase-indexed control states as follows: in the increment
phase, the completed state carrying `true` redirects to the exit phase, while the
completed state carrying `false` redirects to the extend phase in its running state. In
the extend phase, the nested completed state redirects to the exit phase; every other
phase and control state has no redirect.
\end{definition}

\begin{definition}[Successor machine]
\lean{Turing.rt_succTM}
\label{Turing.rt_succTM}
\leanok
\uses{Turing.MultiTapeTM, Turing.phaseHostTM, Turing.rt_succControl, Turing.rt_succParts, Turing.rt_succRedirect}
For a natural number $k$ and a work-tape index $i<k$, define a deterministic Turing
machine with $k$ work tapes, Boolean symbols, and states consisting of a phase together
with its corresponding control data: the increment phase has flag controls, the extend
phase has controls in the disjoint union of seek controls and a nested union of four
indices with seek controls, and the exit phase has a single control. Its initial state
is the running control in the increment phase, and its transitions are obtained by
hosting the phase-specific machine parts and using the phase-dependent redirects.
\end{definition}

\begin{theorem}[Increment phase in a successor computation]
\lean{Turing.rt_succ_increment}
\label{Turing.rt_succ_increment}
\leanok
\uses{Turing.Cfg, Turing.Cfg.mapState, Turing.FlagPhase, Turing.MultiTapeTM.runFrom, Turing.incrementTM, Turing.phaseHostTM_run, Turing.rt_succControl, Turing.rt_succParts, Turing.rt_succRedirect, Turing.rt_succTM}
\difficulty{3}
Let $k$ be a natural number, let $x$ be a Boolean input word, and choose one of the $k$
work tapes. Suppose configurations $a,z$ of the increment machine satisfy that running
from $a$ for $t$ steps gives $z$, that $z$ is in the completed increment phase with flag
$b$, and that no earlier configuration is in any completed increment phase. Starting
from $a$ with the increment phase embedded into the successor machine's state space, the
successor machine agrees through time $t$ with the increment machine, under the same
state embedding. At time $t+1$, its configuration is $z$ with its state replaced by the
exit phase if $b$ is true, and by the extension phase's initial state if $b$ is false.
\end{theorem}

\begin{theorem}[Successor machine run]
\lean{Turing.rt_succ_run}
\label{Turing.rt_succ_run}
\leanok
\uses{Turing.Cfg, Turing.Cfg.mapState, Turing.FinTM.bufferTape, Turing.FinTM.bufferTape_append, Turing.FlagPhase, Turing.MultiTapeTM.runFrom, Turing.MultiTapeTM.runFrom_add, Turing.SeekPhase, Turing.incFixed, Turing.incrementTM, Turing.incrementTM_run_overflow_ofCfg, Turing.incrementTM_run_succ_ofCfg, Turing.phaseHostTM_run, Turing.rt_appendSymbolTM, Turing.rt_appendSymbol_run, Turing.rt_incFixed_length, Turing.rt_succControl, Turing.rt_succParts, Turing.rt_succRedirect, Turing.rt_succTM, Turing.rt_succWord, Turing.rt_succ_increment}
\difficulty{7}
Let $k$ be a natural number, let $x$ and $w$ be Boolean words, and choose a work-tape
index $i$. Suppose a configuration $d$ starts in the successor machine's increment
phase, and tape $i$ agrees with the buffered word $w$ at every offset from $-1$ through
$|w|+1$. Then there is a time $t\le 4|w|+14$ at which the machine first reaches its exit
phase. At that time, tape $i$ contains the successor word of $w$ on the cells beginning
at its original head position, and all other tape cells and configuration data are
unchanged. Up to time $t$, every work-tape head stays within distance $|w|+4$ of its
initial position.
\end{theorem}

\begin{theorem}[Bounded binary-writing machine]
\lean{Turing.rt_exists_boundTM}
\label{Turing.rt_exists_boundTM}
\leanok
\uses{Turing.Cfg, Turing.Cfg.mapState, Turing.FinTM.bufferTape, Turing.MultiTapeTM, Turing.MultiTapeTM.runFrom, Turing.exists_inputLengthTM, Turing.rt_copy_seam, Turing.rt_succControl, Turing.rt_succTM, Turing.rt_succWord_bits, Turing.rt_succ_run, Turing.seamCompTM}
\difficulty{6}
There are a natural number $K$, a finite state set $S$ with decidable equality, a
machine $R$ with $K+1$ work tapes over the Boolean alphabet, designated states
$entry,exit\in S$, and a natural number $C$ such that the following holds for every
Boolean word $x$ and configuration $d$ on input $x$. If $d$ is in state $entry$, its
input head is at position $1$, and all work-tape cells are blank, then $R$ reaches a
configuration equal to $d$ except that its state is $exit$ and tape $0$ stores the
binary expansion of $|x|+1$ starting at that tape's initial head position. The machine
reaches this configuration at some time $t\le C(|x|+1)$, does not visit $exit$ earlier,
and every work-tape head remains within distance $C(\operatorname{bits}(|x|+1)+1)$ of
its initial position through time $t$.
\end{theorem}

\begin{theorem}[Embedded digit-writing machine]
\lean{Turing.rt_exists_digitBoundTM}
\label{Turing.rt_exists_digitBoundTM}
\leanok
\uses{Turing.Cfg, Turing.Cfg.liftSymbols, Turing.FinTM.bufferTape, Turing.MultiTapeTM, Turing.MultiTapeTM.liftSymbols, Turing.MultiTapeTM.liftSymbols_runFrom, Turing.MultiTapeTM.runFrom, Turing.RecSym, Turing.RecSym.digitEmb, Turing.embedEmitCfg, Turing.embedEmitCfg_selected_pos, Turing.embedEmitTM, Turing.embedEmitTM_frame, Turing.embedEmitTM_runFrom, Turing.rt_exists_boundTM}
\difficulty{4}
For any types $Q$ and $T$ with decidable equality, there are a natural number $K$, a
finite state set $S$ with decidable equality, a machine $R$ with $K+1$ work tapes over
symbols encoding digits, marked work cells, states, or tags, designated states
$entry,exit\in S$, and a natural number $C$. For every $k$, embedding
$\iota:\{0,\ldots,K\}\hookrightarrow\{0,\ldots,k-1\}$, Boolean word $x$, and initial
configuration $a$ of $R$ on input $x$, suppose $a$ is in state $entry$, its input head
is at position $1$, all its work-tape cells are blank, and its output is empty. For any
surrounding tapes, head positions, and output prefix, embed this computation into the
$k$-tape machine; then at some time $t\le C(|x|+1)$ it reaches the corresponding
embedded configuration with state $exit$ and the binary expansion of $|x|+1$ written on
tape $\iota(0)$ from its initial head position. It does not visit $exit$ earlier, and
every work-tape head remains within distance $C(\operatorname{bits}(|x|+1)+1)$ of its
initial position through time $t$.
\end{theorem}

\begin{theorem}[Centered Boolean word identity]
\lean{Turing.rt_centeredWord}
\label{Turing.rt_centeredWord}
\leanok
\difficulty{4}
For any type $\Gamma$, function $f:\{\mathrm{false},\mathrm{true}\}\to\Gamma$, and
natural number $s$, mapping $f$ over the list obtained by applying $f$ to the Boolean
test that $j=s$, for $j=0,\ldots,2s$, gives the same list as mapping $f$ over $s$ false
values, followed by true, followed by $s$ false values.
\end{theorem}

\begin{theorem}[Explicit form of the record template]
\lean{Turing.rt_template_fields}
\label{Turing.rt_template_fields}
\leanok
\uses{Turing.binDigits_posWidth_succ, Turing.recordTemplate, Turing.rt_centeredWord}
\difficulty{3}
For any types $Q,T$ and natural numbers $n,s$, the record template consists of two
fields: the binary digits of $n+1$, encoded as digit symbols, followed by a field of
$2s+1$ blank-content cell symbols whose head marker is true exactly at the middle cell.
Equivalently, the second field has marker pattern $s$ false entries, then true, then $s$
false entries.
\end{theorem}

\begin{theorem}[Exact execution of the expanded machine]
\lean{Turing.rt_expand_run}
\label{Turing.rt_expand_run}
\leanok
\uses{Turing.Cfg, Turing.CounterLoopState, Turing.FinTM.bufferTape, Turing.MultiTapeTM.runFrom, Turing.counterLoopCfg, Turing.counterOverhead, Turing.counter_cfg_eq, Turing.expandTM, Turing.expandTM_run_ofCfg}
\difficulty{5}
Let $x$, $pat$, and $ws$ be finite Boolean words, and let $s\in\mathbb{N}$ satisfy
$\operatorname{Nat.ofDigits}_2(ws)=s$. For each $j\in\mathbb{N}$, let $w_j$ be the word
obtained from $ws$ by applying $j$ times the operation that subtracts one from a
nonempty binary word and leaves the word unchanged when subtraction is undefined; let
$C(ws,r)=\sum_{j=0}^{r-1}(2\ell_j+2)$, where $\ell_j$ is the length of the initial
string of false entries in $w_j$. Suppose a configuration $d$ of the three-work-tape
machine constructed from $pat$, on input $x$, has state equal to the counter-decrement
running state, and its third work tape at every offset $p$ from its head with $-1\le
p\le |ws|$ contains the buffer representation of $ws$ at offset $p$. Set
$t=s(|pat|+1)+C(ws,s+1)$. After $t$ steps, the machine is in its done state; the third
work tape contains true throughout the interval of $|ws|$ cells beginning at its initial
head position, the second work tape contains $s$ consecutive copies of $pat$ throughout
the interval of $s|pat|$ cells beginning at its initial head position, and every other
tape cell retains its initial content. The machine is not in the done state at any time
$u<t$, and at every time $u\le t$ each work-tape head is at distance at most
$s|pat|+|ws|+1$ from its initial position.
\end{theorem}

\begin{definition}[Repeated pattern machine]
\lean{Turing.rt_repeatTM}
\label{Turing.rt_repeatTM}
\leanok
\uses{Turing.CounterLoopState, Turing.MultiTapeTM, Turing.SweepPhase, Turing.clearTM, Turing.copyTM, Turing.expandTM, Turing.seamCompTM}
For a finite Boolean word $p$, define a deterministic Turing machine with three work
tapes, Boolean symbols, and states formed from sweep controls, counter-loop controls
indexed by the $|p|+1$ possible counter values, and additional sweep controls. The
machine is the sequential composition that copies tape $0$ onto tape $2$, expands the
pattern $p$ on tape $1$, clears tape $2$, and then runs a counter loop whose body is the
sweep phase and whose counter-decrement phase begins in its running state.
\end{definition}

\begin{theorem}[Repeated word generation by a counter loop]
\lean{Turing.rt_repeat_run}
\label{Turing.rt_repeat_run}
\leanok
\uses{Turing.Cfg, Turing.Cfg.mapState, Turing.CounterLoopState, Turing.FinTM.bufferTape, Turing.MultiTapeTM.runFrom, Turing.SweepPhase, Turing.clearTM, Turing.clearTM_run_ofCfg, Turing.copyTM, Turing.copyTM_run_ofCfg, Turing.counterOverhead, Turing.expandTM, Turing.rt_copy_seam, Turing.rt_expand_run, Turing.rt_repeatTM, Turing.seamCompTM}
\difficulty{6}
Let $\mathrm{pat}$ and $\mathrm{ws}$ be finite Boolean words, let $s\in\mathbb{N}$, and
suppose the binary value of $\mathrm{ws}$, read least significant bit first, is $s$. Let
$d$ be a configuration of a three-work-tape machine on a fixed Boolean input. Assume its
state is the initial sweep phase; tape $0$, at every offset $p$ from its head with
$-1\le p\le |\mathrm{ws}|$, agrees with the buffer containing $\mathrm{ws}$ from cell
$0$ onward; and tape $2$ is blank at all those offsets. Put
$t=4|\mathrm{ws}|+s(|\mathrm{pat}|+1)+H(\mathrm{ws},s+1)+6$, where
$H(\mathrm{ws},r)=\sum_{j=0}^{r-1}(2a_j+2)$ and $a_j$ is the number of initial false
entries in the word obtained from $\mathrm{ws}$ by applying $j$ times the operation that
subtracts one from its fixed-width binary value, with underflow wrapping to the all-true
word. After $t$ steps of the repeat machine for $\mathrm{pat}$, its state is the final
phase and tape $1$ contains $s$ consecutive copies of $\mathrm{pat}$ on the cells from
its initial head position up to, but not including, the cell $s|\mathrm{pat}|$ positions
to its right; its head has advanced by exactly $s|\mathrm{pat}|$. Every other cell and
every other component of the configuration is unchanged. The final phase is not reached
at any earlier time, and at every time $u\le t$, the displacement of each work-tape head
from its initial position is at most $s|\mathrm{pat}|+|\mathrm{ws}|+1$.
\end{theorem}

\begin{theorem}[Run of the released writing machine]
\lean{Turing.rt_writeReleased_run}
\label{Turing.rt_writeReleased_run}
\leanok
\uses{Turing.Cfg, Turing.Cfg.mapState, Turing.FinTM.bufferTape, Turing.MultiTapeTM.runFrom, Turing.MultiTapeTM.visitedByTapeHead, Turing.seamReleaseTM, Turing.seamReleaseTM_firstReturn, Turing.seamReleaseTM_visitedByTapeHead, Turing.writeAdvTM, Turing.writeAdvTM_run_ofCfg}
\difficulty{4}
Let $i$ be a work-tape index, let $w$ be a finite word over an alphabet $\Gamma$, and
let $d$ be a configuration whose state is the release entry state. After exactly $|w|+1$
steps, the released writing machine is in state $0$ and tape $i$ contains $w$ on the
interval of length $|w|$ starting at its initial head position, with the head advanced
by $|w|$ and all other configuration components unchanged. It does not reach that exit
state at any earlier time, and through time $|w|+1$ every work-tape head moves by at
most $|w|$ from its initial position.
\end{theorem}

\begin{theorem}[Adjacent tape writes equal a concatenated write]
\lean{Turing.rt_adjacentWrites}
\label{Turing.rt_adjacentWrites}
\leanok
\uses{Turing.FinTM.bufferTape}
\difficulty{4}
Let $\Gamma$ be a type, let $tape$ be a function from integer cells to optional symbols
in $\Gamma$, let $o$ be an integer, and let $u$ and $v$ be finite words over $\Gamma$.
Replacing the cells from $o$ through the end of $u$ with $u$, then the immediately
following cells with $v$, and leaving all other cells as in $tape$ gives the same tape
as writing the concatenated word $u\mathbin{++}v$ starting at $o$.
\end{theorem}

\begin{definition}[Advance a configuration by a word]
\lean{Turing.rt_advanceCfg}
\label{Turing.rt_advanceCfg}
\leanok
\uses{Turing.Cfg, Turing.FinTM.bufferTape}
Given a configuration of a machine with $k$ work tapes, a work-tape index $i<k$, a
finite word $w$, and a state $q$, define a new configuration with state $q$. On tape
$i$, the cells from its current head position through the next $|w|$ positions are
replaced by the symbols of $w$ in order, and that tape's head is moved forward by $|w|$;
all other tape contents and head positions, as well as the input-head position and
output, are unchanged.
\end{definition}

\begin{theorem}[Appending consecutive tape words]
\lean{Turing.rt_advanceCfg_append}
\label{Turing.rt_advanceCfg_append}
\leanok
\uses{Turing.Cfg, Turing.Cfg.mapState, Turing.rt_adjacentWrites, Turing.rt_advanceCfg}
\difficulty{4}
Let $k$ be a nonnegative integer, let $\Gamma$ be a symbol alphabet, let $P,S$ be state
sets, and fix an input word. For a configuration $d$, tape $i$, words $u,v$, states
$q,r\in P$, and a state map $f:P\to S$, consider the operation that sets the state to
the specified state, writes the word at the current head position and advances that head
by its length. Advancing first by $u$ to state $q$, renaming states by $f$, and then
advancing by $v$ to state $r$ gives the same configuration as first renaming the state
of $d$ by $f$ and then advancing by the concatenation $u\mathbin{+\!\!+}v$ to state $r$.
\end{theorem}

\begin{definition}[Half-state space]
\lean{Turing.rt_halfState}
\label{Turing.rt_halfState}
\leanok
\uses{Turing.CounterLoopState, Turing.SweepPhase}
Define the half-state space to be the disjoint union of sweep-phase states and a
disjoint union whose first part consists of counter-loop states with a two-value counter
and whose second part is another copy of the sweep-phase states.
\end{definition}

\begin{definition}[Window-advance machine]
\lean{Turing.rt_windowAdvTM}
\label{Turing.rt_windowAdvTM}
\leanok
\uses{Turing.MultiTapeTM, Turing.rt_halfState, Turing.rt_repeatTM, Turing.seamCompTM, Turing.seamReleaseTM, Turing.writeAdvTM}
Define a deterministic machine with three work tapes, Boolean symbols, and states built
from the half-state space, a disjoint union of a singleton and two counter values,
another half-state space, and an additional half-state component. It is formed by
composing two repetitions of the one-symbol pattern $[\mathrm{false}]$ around a machine
that writes `true` on tape $1$ and advances, with a release step on tape $0$; the
indicated sweep and counter controls specify the entry and exit states of these
components.
\end{definition}

\begin{theorem}[Bounded run advancing a window]
\lean{Turing.rt_windowAdv_run}
\label{Turing.rt_windowAdv_run}
\leanok
\uses{Turing.Cfg, Turing.Cfg.mapState, Turing.FinTM.bufferTape, Turing.MultiTapeTM.runFrom, Turing.counterOverhead, Turing.counterOverhead_le, Turing.rt_advanceCfg, Turing.rt_advanceCfg_append, Turing.rt_copy_seam, Turing.rt_halfState, Turing.rt_repeatTM, Turing.rt_repeat_run, Turing.rt_windowAdvTM, Turing.rt_writeReleased_run, Turing.seamCompTM, Turing.seamReleaseTM, Turing.writeAdvTM}
\difficulty{6}
Let $ws$ be a Boolean word and $s\in\mathbb N$ such that its binary value is $s$, and
let $d$ be a configuration in the specified entry state. Suppose tape $0$ stores $ws$ at
every offset $p$ with $-1\le p\le |ws|$, and tape $2$ is blank at all those offsets. The
constructed machine reaches its specified exit state at some time $t\le12(s+|ws|)+20$,
never visits that exit state earlier, and up to time $t$ every work-tape head moves by
at most $2s+|ws|+2$ from its initial position; at time $t$, tape $1$ has the word of $s$
false bits, one true bit, and $s$ false bits written from its initial head position,
while every other configuration component agrees with $d$.
\end{theorem}

\begin{definition}[Window machine]
\lean{Turing.rt_windowTM}
\label{Turing.rt_windowTM}
\leanok
\uses{Turing.MultiTapeTM, Turing.SeekPhase, Turing.rt_halfState, Turing.rt_windowAdvTM, Turing.seamCompTM, Turing.skipBackTM, Turing.skipTM}
Define a deterministic machine with three work tapes, Boolean symbols, and states formed
from the window-advance machine's states together with two seek-phase components. It is
the composition of the window-advance machine with a forward skip on tape $1$ followed
by a backward skip on tape $1$, using the specified completed state and running seek
phase as the component boundaries.
\end{definition}

\begin{theorem}[Bounded run constructing a window]
\lean{Turing.rt_window_run}
\label{Turing.rt_window_run}
\leanok
\uses{Turing.Cfg, Turing.Cfg.mapState, Turing.FinTM.bufferTape, Turing.MultiTapeTM.runFrom, Turing.SeekPhase, Turing.rt_advanceCfg, Turing.rt_copy_seam, Turing.rt_halfState, Turing.rt_windowAdvTM, Turing.rt_windowAdv_run, Turing.rt_windowTM, Turing.seamCompTM, Turing.skipBackTM, Turing.skipBackTM_run_ofCfg, Turing.skipTM, Turing.skipTM_run_ofCfg}
\difficulty{6}
Let $ws$ be a Boolean word and $s\in\mathbb N$ such that its binary value is $s$, and
let $d$ be a configuration in the specified entry state. Suppose tape $0$ stores $ws$ at
every offset $p$ with $-1\le p\le |ws|$, tape $2$ is blank at those offsets, and tape
$1$ is blank at every offset $p$ with $-1\le p\le 2s+1$. The constructed machine reaches
its specified exit state at some time $t\le14(s+|ws|)+27$, never visits that exit state
earlier, and up to time $t$ every work-tape head moves by at most $2s+|ws|+3$ from its
initial position; at time $t$, tape $1$ contains the word of $s$ false bits, one true
bit, and $s$ false bits on the interval of length $2s+1$ starting at its initial head
position, and all other configuration components agree with $d$.
\end{theorem}

\begin{definition}[Window cell map]
\lean{Turing.rt_windowCellMap}
\label{Turing.rt_windowCellMap}
\leanok
\uses{Turing.RecSym}
For types $Q$ and $T$, define a map on record symbols that sends a digit symbol carrying
$b\in\{\mathrm{false},\mathrm{true}\}$ to a cell symbol with blank content and mark $b$,
and leaves every other record symbol unchanged.
\end{definition}

\begin{definition}[Cell window machine]
\lean{Turing.rt_cellWindowTM}
\label{Turing.rt_cellWindowTM}
\leanok
\uses{Turing.MultiTapeTM, Turing.MultiTapeTM.liftSymbols, Turing.RecSym, Turing.RecSym.digitEmb, Turing.ScanPhase, Turing.SeekPhase, Turing.rt_halfState, Turing.rt_windowCellMap, Turing.rt_windowTM, Turing.scanKeep, Turing.scanMap, Turing.scanTM, Turing.seamCompTM}
For types $Q$ and $T$ with decidable equality, the cell window machine is the
three-work-tape machine over symbols recording binary digits, simulated Boolean work
cells, states from $Q$, and tags from $T$. Its states combine the phases for the
recursive window procedure, two seek procedures, and a scan phase carrying a unit value;
it is obtained by composing the lifted window machine with a scan that maps cell symbols
and preserves the other tapes.
\end{definition}

\begin{theorem}[Running the cell-window machine]
\lean{Turing.rt_cellWindow_run}
\label{Turing.rt_cellWindow_run}
\leanok
\uses{Turing.Cfg, Turing.Cfg.liftSymbols, Turing.Cfg.liftSymbols_state, Turing.Cfg.liftSymbols_workTapePos, Turing.Cfg.mapState, Turing.FinTM.bufferTape, Turing.FinTM.bufferTape_map, Turing.MultiTapeTM.liftSymbols, Turing.MultiTapeTM.liftSymbols_runFrom, Turing.MultiTapeTM.runFrom, Turing.RecSym.digitEmb, Turing.SeekPhase, Turing.rt_cellWindowTM, Turing.rt_copy_seam, Turing.rt_halfState, Turing.rt_windowCellMap, Turing.rt_windowTM, Turing.rt_window_run, Turing.scanKeep, Turing.scanMap, Turing.scanResult_map, Turing.scanTM, Turing.scanTM_run_ofCfg}
\difficulty{5}
Let $Q$ and $T$ be types with decidable equality, let $x$ be a Boolean input word, let
$w$ be a Boolean word, and let $s$ be a natural number such that the binary value of $w$
is $s$. Suppose a configuration $a$ of the specified three-tape machine is in its sweep
entry state; within the interval of offsets from $-1$ to $|w|$, tape 0 stores $w$
starting at its head, tape 2 is blank, and tape 1 is blank at every offset from $-1$ to
$2s+1$. After encoding the digits as symbols and adjoining the machine's state tag,
there is a time $t\le 18(s+|w|)+32$ when the machine reaches its designated exit state,
having avoided that state at every earlier time; up to and including time $t$, every
work-tape head stays within distance $2s+|w|+3$ of its initial position. At time $t$,
all configuration data except the state and work tapes agree with the encoded starting
configuration, and tape 1 has been overwritten on the $2s+1$ cells beginning at its
initial head position by $s$ blank-cell symbols, one marked cell symbol, and then $s$
blank-cell symbols; tape contents outside that interval are unchanged.
\end{theorem}

\begin{theorem}[Head-radius bound under tape embedding]
\lean{Turing.rt_embed_radius}
\label{Turing.rt_embed_radius}
\leanok
\uses{Turing.Cfg, Turing.MultiTapeTM, Turing.MultiTapeTM.runFrom, Turing.embedEmitCfg, Turing.embedEmitCfg_selected_pos, Turing.embedEmitTM, Turing.embedEmitTM_frame, Turing.embedEmitTM_runFrom}
\difficulty{4}
Let $\Gamma,S$ be types, let $m,k$ be natural numbers, and let $x$ be an input word over
$\Gamma$. Let $\iota$ embed the $m$ tape indices into the $k$ tape indices, let $M$ be a
machine with $m$ work tapes, and choose contents and head positions for the additional
tapes and an output prefix. For a starting configuration $a$, time bound $t$, and
nonnegative integer $B$, suppose that at every time $u\le t$ each work-tape head of $M$
is within distance $B$ of its initial position. Then, starting from the configuration
formed by embedding $a$ with the chosen additional tapes and prefix, every work-tape
head of the embedded machine remains within distance $B$ of its initial position at
every time $u\le t$.
\end{theorem}

\begin{theorem}[Embedding the cell-window computation]
\lean{Turing.rt_cellWindow_embed_run}
\label{Turing.rt_cellWindow_embed_run}
\leanok
\uses{Turing.Cfg, Turing.Cfg.liftSymbols, Turing.Cfg.mapState, Turing.FinTM.bufferTape, Turing.MultiTapeTM.runFrom, Turing.RecSym, Turing.RecSym.digitEmb, Turing.SeekPhase, Turing.embedEmitCfg, Turing.embedEmitTM, Turing.embedEmitTM_runFrom, Turing.rt_cellWindowTM, Turing.rt_cellWindow_run, Turing.rt_embed_radius, Turing.rt_halfState}
\difficulty{4}
Let $Q,T$ be types with decidable equality, let $k\in\mathbb{N}$, and let $\iota$ embed
the three tape indices into the $k$ tape indices. Let $ws$ be a Boolean word and
$s\in\mathbb{N}$ such that its binary value is $s$. Suppose a three-tape configuration
$a$ on input $x$ starts in the initial sweep phase; tape 0 stores $ws$ as a buffer on
the interval from one cell left of its head through the cell indexed by $|ws|$, tape 2
is blank on that interval, tape 1 is blank from one cell left through the cell indexed
by $2s+1$, and the output is empty. For any surrounding tapes, head positions, and
output prefix used to embed the configuration into $k$ tapes, there is a time
$t\le18(s+|ws|)+32$ such that the embedded machine has not reached its designated done
state before $t$, every work-tape head stays within distance $2s+|ws|+3$ of its initial
position through time $t$, and at time $t$ the designated state is reached and the
output remains the prefix. The embedded configuration agrees with its initial version
except on the interval of tape $\iota(1)$ from its head position through the next $2s+1$
cells, where the buffer now stores $s$ false cell symbols, then one true cell symbol,
then $s$ false cell symbols.
\end{theorem}

\begin{theorem}[Equality of an embedded configuration]
\lean{Turing.rt_embed_matches}
\label{Turing.rt_embed_matches}
\leanok
\uses{Turing.Cfg, Turing.MultiTapeTM, Turing.embedEmitCfg, Turing.embedEmitCfg_selected_pos, Turing.embedEmitCfg_selected_tape, Turing.embedEmitTM_frame}
\difficulty{4}
Let $\Gamma,S$ be types, let $m,k$ be natural numbers, and let $x$ be an input word over
$\Gamma$. Let $\iota$ be an injection from the $m$ tape indices into the $k$ tape
indices, and let $a,d$ be configurations with $m$ and $k$ work tapes, respectively.
Choose contents and head positions for the $k$ tapes, and a prefix word. If $a$ and $d$
have the same state and input-head position; each tape of $a$ agrees with the
corresponding embedded tape of $d$, including its head position; every tape of $d$
outside the image of $\iota$ agrees with the chosen contents and head positions; and the
prefix followed by the output of $a$ equals the output of $d$, then embedding $a$ into
the $k$-tape configuration with those outside tapes and prefix produces exactly $d$.
\end{theorem}

\begin{theorem}[Embedded clean-tape digit writer]
\lean{Turing.rt_exists_cleanDigitBoundTM}
\label{Turing.rt_exists_cleanDigitBoundTM}
\leanok
\uses{Turing.Cfg, Turing.Cfg.liftSymbols, Turing.FinTM.bufferTape, Turing.FinTM.bufferTape_map, Turing.MultiTapeTM, Turing.MultiTapeTM.runFrom, Turing.RecSym, Turing.RecSym.digitEmb, Turing.embedEmitCfg, Turing.embedEmitTM, Turing.rt_embed_matches, Turing.rt_exists_digitBoundTM}
\difficulty{5}
For any types $Q$ and $T$ with decidable equality, there are a natural number $K$, a
finite state set $S$ with decidable equality, a machine $R$ with $K+1$ work tapes over
symbols encoding digits, marked work cells, states, or tags, designated states
$entry,exit\in S$, and a natural number $C$. For every $k$, embedding
$\iota:\{0,\ldots,K\}\hookrightarrow\{0,\ldots,k-1\}$, Boolean word $x$, and
configuration $d$ on the encoded input, if $d$ begins in state $entry$ with input head
at position $1$ and every cell of every tape in the image of $\iota$ blank, then the
machine obtained by placing $R$ on those embedded tapes reaches a configuration equal to
$d$ except that its state is $exit$ and tape $\iota(0)$ stores the binary expansion of
$|x|+1$ from its initial head position. It does so by some time $t\le C(|x|+1)$, does
not visit $exit$ earlier, and all work-tape heads remain within distance
$C(\operatorname{bits}(|x|+1)+1)$ of their initial positions through time $t$.
\end{theorem}

\begin{definition}[Move and clear a word]
\lean{Turing.rt_moveWordTM}
\label{Turing.rt_moveWordTM}
\leanok
\uses{Turing.MultiTapeTM, Turing.SweepPhase, Turing.clearTM, Turing.copyTM, Turing.seamCompTM}
Given an alphabet, a natural number $k$ of work tapes, and source and destination tape
indices among those $k$ tapes, define a $k$-tape machine that copies the word on the
source tape to the destination tape and then clears the source tape. Its states track
the copying phase and the subsequent clearing phase.
\end{definition}

\begin{theorem}[Copying and clearing a tape word]
\lean{Turing.rt_moveWord_run}
\label{Turing.rt_moveWord_run}
\leanok
\uses{Turing.Cfg, Turing.Cfg.mapState, Turing.FinTM.bufferTape, Turing.MultiTapeTM.runFrom, Turing.SweepPhase, Turing.clearTM, Turing.clearTM_run_ofCfg, Turing.copyTM, Turing.copyTM_run_ofCfg, Turing.rt_copy_seam, Turing.rt_moveWordTM}
\difficulty{5}
Let $k$ be a natural number, let $\Gamma$ be an alphabet, and let $x$ be an input word.
Choose distinct work tapes $\mathrm{src},\mathrm{dst}$ and a word $w\in\Gamma^*$.
Suppose a configuration $d$ starts in the copy phase and, for every integer offset $p$
with $-1\le p\le |w|$, tape $\mathrm{src}$ at its head position plus $p$ contains the
symbol of $w$ at offset $p$ (and is blank outside the word). After $4|w|+5$ steps, the
machine is in its final phase, tape $\mathrm{src}$ is blank on the $|w|$ cells beginning
at its original head position, and tape $\mathrm{dst}$ contains a copy of $w$ on the
$|w|$ cells beginning at its original head position; all other cells and configuration
data are unchanged. The final phase is not reached earlier, and at every time up to and
including this one, every work-tape head is within distance $|w|+1$ of its initial
position.
\end{theorem}

\begin{definition}[Cell window state space]
\lean{Turing.rt_cellWindowState}
\label{Turing.rt_cellWindowState}
\leanok
\uses{Turing.ScanPhase, Turing.SeekPhase, Turing.rt_halfState}
For the types and phase spaces used by the window procedure, define its cell window
state space as the disjoint union of the recursive half-window states, a unit state or
one of two indexed states, another copy of the half-window states, two seek-phase state
spaces, and a scan phase carrying a unit value.
\end{definition}

\begin{theorem}[Ready-state cell-window computation]
\lean{Turing.rt_cellWindow_ready_run}
\label{Turing.rt_cellWindow_ready_run}
\leanok
\uses{Turing.Cfg, Turing.Cfg.liftSymbols, Turing.Cfg.mapState, Turing.FinTM.bufferTape, Turing.FinTM.bufferTape_map, Turing.MultiTapeTM.runFrom, Turing.RecSym, Turing.RecSym.digitEmb, Turing.SeekPhase, Turing.embedEmitCfg, Turing.embedEmitTM, Turing.rt_cellWindowState, Turing.rt_cellWindowTM, Turing.rt_cellWindow_embed_run, Turing.rt_embed_matches, Turing.rt_halfState}
\difficulty{5}
Let $Q,T$ have decidable equality, let $k\in\mathbb{N}$, and let
$\iota:\mathrm{Fin}(3)\hookrightarrow\mathrm{Fin}(k)$ be an embedding. Let $ws$ be a
Boolean word and $s\in\mathbb{N}$ whose binary value is $s$. Suppose a $k$-tape
configuration $d$ on the input obtained by embedding each Boolean symbol into the
digit-symbol alphabet starts in the specified initial sweep state. On the three embedded
tapes, tape 0 stores the digit encoding of $ws$ as a buffer relative to its head, while
tapes 1 and 2 are blank everywhere. Then there is a time $t\le18(s+|ws|)+32$ such that
the designated done state is not reached before time $t$, every head stays within
distance $2s+|ws|+3$ of its initial position through time $t$, and at time $t$ the
designated state is reached. At that time, the configuration is unchanged except that
tape $\iota(1)$ contains, from its initial head position through the next $2s$ cells, a
buffer of $s$ false cell symbols, one true cell symbol, and $s$ false cell symbols.
\end{theorem}

\begin{theorem}[Buffer contents outside the written interval]
\lean{Turing.rt_bufferMask}
\label{Turing.rt_bufferMask}
\leanok
\uses{Turing.FinTM.bufferTape}
\difficulty{3}
For any symbol alphabet $\Gamma$, word $w$, and integers $o,p$, the buffer-tape content
at offset $p-o$ is unchanged by restricting it to the interval $[o,o+|w|)$: if $p$ is
outside this interval, that content is blank.
\end{theorem}

\begin{definition}[Rich window state space]
\lean{Turing.rt_richWindowState}
\label{Turing.rt_richWindowState}
\leanok
\uses{Turing.SweepPhase, Turing.rt_cellWindowState}
Define the rich window state space as the disjoint union of a sweep phase, the cell
window state space, two sweep-phase state spaces, and one further sweep phase.
\end{definition}

\begin{definition}[Rich window machine]
\lean{Turing.rt_richWindowTM}
\label{Turing.rt_richWindowTM}
\leanok
\uses{Turing.MultiTapeTM, Turing.RecSym, Turing.clearTM, Turing.copyTM, Turing.embedEmitTM, Turing.rt_cellWindowTM, Turing.rt_moveWordTM, Turing.rt_richWindowState, Turing.seamCompTM}
Let $Q$ and $T$ be types with decidable equality, let $k$ be a natural number, choose a
radius tape and a destination tape among the $k$ work tapes, and choose an injective
assignment of three tape indices to those work tapes. The rich window machine over
symbols recording binary digits, simulated Boolean cells, states from $Q$, and tags from
$T$ is formed by composing a copy of the radius tape into the first assigned tape, an
embedded cell window computation, a move-and-clear operation from the second assigned
tape to the destination, and a final clearing of the radius tape. Its states are the
rich window state space.
\end{definition}

\begin{theorem}[Constructing a rich window]
\lean{Turing.rt_richWindow_run}
\label{Turing.rt_richWindow_run}
\leanok
\uses{Turing.Cfg, Turing.Cfg.mapState, Turing.FinTM.bufferTape, Turing.MultiTapeTM.runFrom, Turing.RecSym, Turing.RecSym.digitEmb, Turing.SweepPhase, Turing.clearTM, Turing.clearTM_run_ofCfg, Turing.copyTM, Turing.copyTM_run_ofCfg, Turing.embedEmitTM, Turing.rt_bufferMask, Turing.rt_cellWindowState, Turing.rt_cellWindowTM, Turing.rt_cellWindow_ready_run, Turing.rt_copy_seam, Turing.rt_moveWordTM, Turing.rt_moveWord_run, Turing.rt_richWindowState, Turing.rt_richWindowTM, Turing.seamCompTM}
\difficulty{6}
Let $Q,T$ be types with decidable equality, let $k$ be a natural number, and choose
work-tape indices $\mathrm{src},\mathrm{dst}$ and an embedding
$\iota:\{0,1,2\}\hookrightarrow\{0,\ldots,k-1\}$, with neither chosen index in the image
of $\iota$. Let $\mathrm{ws}$ be a Boolean word and $s$ a natural number whose value is
the binary value of $\mathrm{ws}$ read least-significant bit first. Suppose a
configuration $d$ starts in the copy phase, tape $\mathrm{src}$ contains the buffered
word $\mathrm{ws}$ at offsets $-1$ through $|\mathrm{ws}|$ (encoded as digit symbols),
and all tapes indexed by $\iota$ are blank. Then there is a time $t\le
26(s+|\mathrm{ws}|)+48$ at which the machine reaches its final phase for the first time.
At that time, tape $\mathrm{dst}$ contains $s$ zeros, then a one, then $s$ zeros, each
encoded as a cell symbol with blank content, on the $2s+1$ cells beginning at its
original head position; all other cells and configuration data are unchanged. Up to time
$t$, every work-tape head stays within distance $2s+|\mathrm{ws}|+3$ of its initial
position.
\end{theorem}

\begin{theorem}[Bounded field-writing machine]
\lean{Turing.rt_exists_boundFieldTM}
\label{Turing.rt_exists_boundFieldTM}
\leanok
\uses{Turing.Cfg, Turing.Cfg.mapState, Turing.FinTM.bufferTape, Turing.MultiTapeTM, Turing.MultiTapeTM.runFrom, Turing.RecSym, Turing.RecSym.digitEmb, Turing.SweepPhase, Turing.embedEmitTM, Turing.rt_bufferMask, Turing.rt_copy_seam, Turing.rt_exists_cleanDigitBoundTM, Turing.rt_moveWordTM, Turing.rt_moveWord_run, Turing.seamCompTM}
\difficulty{5}
For any types $Q$ and $T$ with decidable equality, there are a natural number $K$, a
finite state set $S$ with decidable equality, a machine $R$ with $K+1$ work tapes over
symbols that encode either Boolean digits, marked Boolean work cells, states from $Q$,
or tags from $T$, designated states $entry,exit\in S$, and a natural number $C$, with
the following property. For every $k$, every embedding
$\iota:\{0,\ldots,K\}\hookrightarrow\{0,\ldots,k-1\}$, every tape index $dst$ outside
the image of $\iota$, every Boolean word $x$, and every configuration $d$ on input $x$
encoded as digit symbols, if $d$ is in state $entry$, its input head is at position $1$,
and all cells of the work tapes indexed by $\iota$ are blank, then the embedded machine
reaches a configuration equal to $d$ except that its state is $exit$ and tape $dst$
stores the binary expansion of $|x|+1$ starting at its initial head position (all other
tape cells retain their contents). It does so at some time $t\le C(|x|+1)$, does not
visit state $exit$ earlier, and throughout the first $t$ steps every work-tape head
stays within distance $C(\operatorname{bits}(|x|+1)+1)$ of its initial position.
\end{theorem}

\begin{theorem}[Machine for Writing a Record Template]
\lean{Turing.exists_recordTemplateTM}
\label{Turing.exists_recordTemplateTM}
\leanok
\uses{Turing.Cfg, Turing.Cfg.mapState, Turing.FinTM.bufferTape, Turing.MultiTapeTM, Turing.MultiTapeTM.runFrom, Turing.RecSym, Turing.RecSym.digitEmb, Turing.SeekPhase, Turing.fieldOffset, Turing.fieldsTape, Turing.fieldsTape_field, Turing.joinFields, Turing.joinFields_recordTemplate_length, Turing.recordTemplate, Turing.rt_copy_seam, Turing.rt_exists_boundFieldTM, Turing.rt_richWindowState, Turing.rt_richWindowTM, Turing.rt_richWindow_run, Turing.rt_template_fields, Turing.seamCompTM, Turing.skipBackTM, Turing.skipBackTM_run_ofCfg, Turing.skipTM, Turing.skipTM_run_ofCfg}
\difficulty{5}
For finite decidable types $Q,T$, there are a natural number $K$, a finite decidable
state type $P$, a deterministic machine $R$ with $K+2$ work tapes over the record
alphabet, entry and exit states, and a constant $C$ such that the following holds. For
every binary string $x$, natural number $s$, and binary string $w$ representing $s$ in
base two, any configuration $d$ on the embedded input with entry state and input head at
position $1$, with tape $1$ containing $w$, tape $0$ blank over the joined template
interval, and every tape of index at least $2$ blank, reaches in at most $C(|\langle
H\rangle|+|w|+|x|+1)$ steps the configuration obtained by entering the exit state and
writing the joined fields of the prescribed record template $H$ for parameters
$Q,T,|x|,s$ on tape $0$. It avoids the exit state before that time, and all work-tape
heads stay within distance $C(|\langle H\rangle|+|w|+1)$ of their starting positions
through that time.
\end{theorem}

\begin{theorem}[Bounded scratch-template construction]
\lean{Turing.rt_exists_scratchTemplateTM}
\label{Turing.rt_exists_scratchTemplateTM}
\leanok
\uses{Turing.Cfg, Turing.Cfg.select, Turing.FinTM.bufferTape, Turing.MultiTapeTM, Turing.MultiTapeTM.runFrom, Turing.RecSym, Turing.RecSym.digitEmb, Turing.embedEmitCfg, Turing.embedEmitCfg_eq, Turing.embedEmitCfg_select, Turing.embedEmitTM, Turing.embedEmitTM_runFrom, Turing.exists_recordTemplateTM, Turing.fieldsTape, Turing.joinFields, Turing.recordTemplate, Turing.rt_embed_radius}
\difficulty{4}
For finite types $Q$ and $T$ with decidable equality, there are a natural number $K$, a
finite state set $P$ with decidable equality, a machine $R$ with $K+2$ work tapes over
symbols encoding digits, marked work cells, states, or tags, designated states
$entry,exit\in P$, and a natural number $C$. For every $k$, embedding
$\iota:\{0,\ldots,K+1\}\hookrightarrow\{0,\ldots,k-1\}$, Boolean word $x$, natural
number $s$, and Boolean word $ws$ whose binary value is $s$, consider a configuration
$d$ on the encoded input. If it begins in state $entry$ with input head at position $1$,
tape $\iota(1)$ holds $ws$ in the prescribed buffer format on positions from $-1$
through $|ws|$ relative to its head, tape $\iota(0)$ is blank on positions from $-1$
through the length of the joined record template for $x$ and $s$, and all embedded tapes
with index at least $2$ are blank, then the embedded machine reaches a configuration
equal to $d$ except that its state is $exit$ and tape $\iota(0)$ contains the joined
record template on that interval. It does so by some time at most $C$ times the sum of
the template length, $|ws|$, $|x|$, and $1$; it does not visit $exit$ earlier, and every
work-tape head stays within distance $C$ times the sum of the template length, $|ws|$,
and $1$ of its initial position throughout.
\end{theorem}

\begin{theorem}[Binary digits of one]
\lean{Turing.rt_binDigits_one}
\label{Turing.rt_binDigits_one}
\leanok
\uses{Turing.binDigits}
\difficulty{4}
For every positive integer $w$, the list of the first $w$ binary digits of $1$, starting
with the least significant digit, is a true digit followed by $w-1$ false digits.
\end{theorem}

\begin{theorem}[Record of an initial configuration]
\lean{Turing.rt_initRecord_fields}
\label{Turing.rt_initRecord_fields}
\leanok
\uses{Turing.Cfg, Turing.Cfg.init, Turing.cfgRecord, Turing.cfgRecord_init, Turing.posWidth, Turing.rt_binDigits_one, Turing.rt_centeredWord}
\difficulty{3}
Let $x$ be a Boolean input word, let $s$ be a nonnegative integer, let $q_0$ be a state,
and let $\tau_0$ be a tag. The record of the initial configuration on $x$, with tag
$\tau_0$, consists of a state field containing $q_0$, a position field containing the
binary digits of $1$ in width equal to the bit width of $|x|+1$, and one field for each
of the $k$ work tapes. Each work-tape field has $2s+1$ blank cells with its head mark at
the central cell; the final field contains the tag $\tau_0$.
\end{theorem}

\begin{definition}[Recursive tail state space]
\lean{Turing.rt_tailState}
\label{Turing.rt_tailState}
\leanok
\uses{Turing.SeekPhase, Turing.SweepPhase, Turing.WritePhase}
Define a family of state spaces indexed by natural numbers recursively: at index $0$,
use the write phase for a word of length parameter $1$; at index $k+1$, take the
disjoint union of a sweep phase, a seek phase, the state space at index $k$, and another
seek phase.
\end{definition}

\begin{definition}[Recursive tail entry state]
\lean{Turing.rt_tailEntry}
\label{Turing.rt_tailEntry}
\leanok
\uses{Turing.rt_tailState}
For each natural number $k$, define an entry state in the recursive tail state space: at
$k=0$ it is the write phase at position $0$, and at every positive index it is the sweep
phase in the first component.
\end{definition}

\begin{definition}[Recursive tail exit state]
\lean{Turing.rt_tailExit}
\label{Turing.rt_tailExit}
\leanok
\uses{Turing.rt_tailState}
For each natural number $k$, define an exit state in the recursive tail state space: at
$k=0$ it is the done state of the write phase, and at every positive index it is the
done state in the final seek-phase component.
\end{definition}

\begin{definition}[Recursive tail machine]
\lean{Turing.rt_tailTM}
\label{Turing.rt_tailTM}
\leanok
\uses{Turing.MultiTapeTM, Turing.copyTM, Turing.rt_tailEntry, Turing.rt_tailExit, Turing.rt_tailState, Turing.seamCompTM, Turing.skipBackTM, Turing.skipTM, Turing.writeConstTM}
Given an alphabet $\Gamma$, a number $m$ of work tapes, indices
$\mathrm{src},\mathrm{dst}\in\mathrm{Fin}(m)$, and a symbol $\mathrm{last}\in\Gamma$,
define a deterministic Turing machine for each $k\in\mathbb N$. For $k=0$ it writes the
one-symbol word $[\mathrm{last}]$ on the destination tape; for $k+1$, it copies from the
source tape to the destination tape, scans forward on the destination tape, runs the
machine defined for $k$, and then scans back on the destination tape.
\end{definition}

\begin{theorem}[Bounded construction of a tail field]
\lean{Turing.rt_tail_run}
\label{Turing.rt_tail_run}
\leanok
\uses{Turing.Cfg, Turing.Cfg.mapState, Turing.FinTM.bufferTape, Turing.MultiTapeTM.runFrom, Turing.SeekPhase, Turing.SweepPhase, Turing.copyTM, Turing.copyTM_run_ofCfg, Turing.fieldOffset, Turing.fieldsTape, Turing.fieldsTape_cons_shift, Turing.fieldsTape_field, Turing.joinFields, Turing.joinFields_cons_cons, Turing.rt_copy_seam, Turing.rt_tailEntry, Turing.rt_tailExit, Turing.rt_tailState, Turing.rt_tailTM, Turing.seamCompTM, Turing.skipBackTM, Turing.skipBackTM_run_ofCfg, Turing.skipTM, Turing.skipTM_run_ofCfg, Turing.writeConstTM_run_ofCfg}
\difficulty{7}
Let $m,k\in\mathbb{N}$, let $\Gamma$ be a type, let $src,dst\in\mathrm{Fin}(m)$ be
distinct, let $last\in\Gamma$, and let $w$ be a word over $\Gamma$. Define
$fs=\mathrm{replicate}(k,w)++[[last]]$ and $F=\mathrm{joinFields}(fs)$. Let $d$ be a configuration
on a fixed input in the prescribed entry state. Suppose the source tape contains the
buffer for $w$ at every offset from $-1$ through $|w|$, and the destination tape is blank
at every offset from $-1$ through $|F|$. Then for some $t\le(k+1)(4|w|+12)$, the
machine reaches its prescribed exit state at time $t$ and is not in that state earlier.
Throughout times $u\le t$, every work-tape head is within distance $|F|+1$ of its initial
position. At time $t$, the configuration agrees with $d$ except that its state is the
exit state and the destination tape on the interval from its initial head position up to
but not including that position plus $|F|$ agrees with $\mathrm{fieldsTape}(fs,\cdot)$; all
other cells and tapes are unchanged.
\end{theorem}

\begin{definition}[Positioning machine]
\lean{Turing.rt_positionTM}
\label{Turing.rt_positionTM}
\leanok
\uses{Turing.MultiTapeTM, Turing.ScanPhase, Turing.SweepPhase, Turing.WritePhase, Turing.copyTM, Turing.scanKeep, Turing.scanMap, Turing.scanTM, Turing.seamCompTM, Turing.writeConstTM}
For a symbol alphabet $\Gamma$, a number $m\in\mathbb N$ of work tapes, source and
destination indices $\mathrm{src},\mathrm{dst}\in\mathrm{Fin}(m)$, and symbols
$\mathrm{zero},\mathrm{one}\in\Gamma$, define an $m$-work-tape Turing machine by
composing a copy pass from the source tape to the destination, a scan of the destination
using the constant map with value $\mathrm{zero}$ and the unit state, and a machine that
writes the one-symbol word $[\mathrm{one}]$ on the destination tape.
\end{definition}

\begin{theorem}[Writing a marked word on a tape]
\lean{Turing.rt_position_run}
\label{Turing.rt_position_run}
\leanok
\uses{Turing.Cfg, Turing.Cfg.mapState, Turing.FinTM.bufferTape, Turing.MultiTapeTM.runFrom, Turing.ScanPhase, Turing.SweepPhase, Turing.WritePhase, Turing.copyTM, Turing.copyTM_run_ofCfg, Turing.rt_copy_seam, Turing.rt_positionTM, Turing.scanKeep, Turing.scanMap, Turing.scanResult_map, Turing.scanTM, Turing.scanTM_run_ofCfg, Turing.seamCompTM, Turing.writeConstTM, Turing.writeConstTM_run_ofCfg}
\difficulty{6}
Let $m$ be a natural number, let $\Gamma$ be an alphabet, and let $x$ be an input word.
Choose distinct work tapes $\mathrm{src},\mathrm{dst}$, symbols
$\mathrm{zero},\mathrm{one}\in\Gamma$, and a nonempty word $w\in\Gamma^*$. Suppose a
configuration $d$ starts in the copy phase, tape $\mathrm{src}$ agrees with the buffered
word $w$ at every offset from $-1$ through $|w|$, and tape $\mathrm{dst}$ is blank at
those offsets. After $4|w|+10$ steps, the machine is in its final phase and tape
$\mathrm{dst}$ contains $\mathrm{one}$ followed by $|w|-1$ copies of $\mathrm{zero}$ on
the $|w|$ cells beginning at its original head position; all other cells and
configuration data are unchanged. This phase is not reached earlier, and every work-tape
head stays within distance $|w|+1$ of its initial position at all times up to and
including that time.
\end{theorem}

\begin{definition}[Clean template machine]
\lean{Turing.rt_cleanTemplateTM}
\label{Turing.rt_cleanTemplateTM}
\leanok
\uses{Turing.MultiTapeTM, Turing.SeekPhase, Turing.SweepPhase, Turing.clearTM, Turing.seamCompTM, Turing.skipBackTM}
For a symbol alphabet $\Gamma$, a number $m\in\mathbb N$ of work tapes, and a source
index $\mathrm{src}\in\mathrm{Fin}(m)$, define an $m$-work-tape Turing machine by
composing a clearing pass on the source tape, a backward skip on that tape, and a second
clearing pass on it.
\end{definition}

\begin{theorem}[Running the clean-template machine]
\lean{Turing.rt_cleanTemplate_run}
\label{Turing.rt_cleanTemplate_run}
\leanok
\uses{Turing.Cfg, Turing.Cfg.mapState, Turing.FinTM.bufferTape, Turing.MultiTapeTM.runFrom, Turing.SeekPhase, Turing.SweepPhase, Turing.clearTM, Turing.clearTM_run_ofCfg, Turing.fieldOffset, Turing.fieldsTape, Turing.fieldsTape_field, Turing.joinFields, Turing.rt_cleanTemplateTM, Turing.rt_copy_seam, Turing.seamCompTM, Turing.skipBackTM, Turing.skipBackTM_run_ofCfg}
\difficulty{6}
Let $\Gamma$ be a type, let $m$ be a natural number, let $x$ be an input word over
$\Gamma$, and choose a source tape among the $m$ work tapes. Suppose the starting
configuration is in the sweep entry state and that, at every offset $q$ from $-1$
through the length of the joined fields $u,v$, the source tape stores the joined-field
tape content at offset $q-(h-|u|-1)$, where $h$ is the source head's initial position.
After exactly $3|u|+2|v|+9$ steps, the machine is in its designated exit state, the
source tape is blank everywhere, and its head is at position $h-|u|-1$; every earlier
configuration avoids the exit state, and at every time up to and including the final
time every work-tape head is within distance $|u|+|v|+2$ of its initial position.
\end{theorem}

\begin{definition}[Initial remainder state space]
\lean{Turing.rt_initRestState}
\label{Turing.rt_initRestState}
\leanok
\uses{Turing.ScanPhase, Turing.SeekPhase, Turing.SweepPhase, Turing.WritePhase, Turing.rt_tailState}
For each natural number $k$, define the initial remainder state space as the disjoint
union of two one-symbol writing-phase state spaces, six seeking-phase state spaces,
three sweep-phase state spaces, a scan-phase state space indexed by the unit type, and
the tail state space at parameter $k$.
\end{definition}

\begin{definition}[Initialization and remainder exit state]
\lean{Turing.rt_initRestExit}
\label{Turing.rt_initRestExit}
\leanok
\uses{Turing.rt_initRestState}
For each $k\in\mathbb N$, define the exit state in the initialization-and-remainder
state space to be the done state in its final sweep-phase component.
\end{definition}

\begin{definition}[Initialization and remainder machine]
\lean{Turing.rt_initRestTM}
\label{Turing.rt_initRestTM}
\leanok
\uses{Turing.MultiTapeTM, Turing.rt_cleanTemplateTM, Turing.rt_initRestState, Turing.rt_positionTM, Turing.rt_tailEntry, Turing.rt_tailExit, Turing.rt_tailTM, Turing.seamCompTM, Turing.skipBackTM, Turing.skipTM, Turing.writeConstTM}
For a symbol alphabet $\Gamma$, a number $m\in\mathbb N$ of work tapes, source and
destination indices $\mathrm{src},\mathrm{dst}\in\mathrm{Fin}(m)$, symbols
$\mathrm{zero},\mathrm{one},\mathrm{first},\mathrm{last}\in\Gamma$, and $k\in\mathbb N$,
define an $m$-work-tape Turing machine with the initialization-and-remainder state
space. It is assembled by composing a machine that writes $[\mathrm{first}]$ on the
destination, a forward skip there, the positioning machine, forward skips on the source
and destination, the recursive tail machine, and backward skips on the destination
interleaved with the clean-template machine.
\end{definition}

\begin{theorem}[Initialization of the remaining record fields]
\lean{Turing.rt_initRest_run}
\label{Turing.rt_initRest_run}
\leanok
\uses{Turing.Cfg, Turing.Cfg.mapState, Turing.FinTM.bufferTape, Turing.MultiTapeTM, Turing.MultiTapeTM.runFrom, Turing.ScanPhase, Turing.SeekPhase, Turing.SweepPhase, Turing.WritePhase, Turing.fieldOffset, Turing.fieldsTape, Turing.fieldsTape_field, Turing.joinFields, Turing.joinFields_cons_cons, Turing.joinFields_length, Turing.rt_cleanTemplateTM, Turing.rt_cleanTemplate_run, Turing.rt_copy_seam, Turing.rt_fields_slice, Turing.rt_initRestExit, Turing.rt_initRestState, Turing.rt_initRestTM, Turing.rt_positionTM, Turing.rt_position_run, Turing.rt_tailEntry, Turing.rt_tailExit, Turing.rt_tailState, Turing.rt_tailTM, Turing.rt_tail_run, Turing.seamCompTM, Turing.skipBackTM, Turing.skipBackTM_run_ofCfg, Turing.skipTM, Turing.skipTM_run_ofCfg, Turing.writeConstTM, Turing.writeConstTM_run_ofCfg}
\difficulty{7}
Let $m>0$, let $x$ be a Boolean input word, and choose distinct tape indices
$\mathrm{src},\mathrm{dst}<m$, symbols
$\mathrm{zero},\mathrm{one},\mathrm{first},\mathrm{last}\in\Gamma$, and words $u,v$ over
$\Gamma$ with $|u|>0$. Let $k$ be a nonnegative integer and let $d$ be an $m$-tape
configuration in the initial write state. Suppose the source tape, relative to its head,
encodes the two fields $u,v$ joined by a blank separator, and the destination tape is
blank throughout the specified range from offset $-1$ through the length of the joined
output fields. Form the output fields from the one-symbol field $[\mathrm{first}]$, the
field beginning with $\mathrm{one}$ followed by $|u|-1$ copies of $\mathrm{zero}$, then
$k$ copies of $v$, and finally the one-symbol field $[\mathrm{last}]$, joining
consecutive fields with blank separators. The initialization machine reaches its
designated exit state at some time $t\le 16|J|+8|v|+16$, where $J$ is the joined output
word; it does not reach that exit state at any earlier time, and throughout times $z\le
t$ every work-tape head moves by at most $2|J|+2|v|+8$ from its initial position. At
time $t$, the source tape is blank everywhere; the destination tape contains the
encoding of $J$ on the interval from its initial head position through that position
plus $|J|-1$; and every other tape cell retains its initial content.
\end{theorem}

\begin{theorem}[Machine for Initializing Records]
\lean{Turing.exists_recordInitTM}
\label{Turing.exists_recordInitTM}
\leanok
\uses{Turing.Cfg, Turing.Cfg.init, Turing.Cfg.mapState, Turing.FinTM.bufferTape, Turing.MultiTapeTM, Turing.MultiTapeTM.runFrom, Turing.RecSym, Turing.RecSym.digitEmb, Turing.cfgRecord, Turing.embedEmitTM, Turing.fieldsTape, Turing.joinFields, Turing.joinFields_cfgRecord_length, Turing.joinFields_recordTemplate_length, Turing.recordLength, Turing.recordTemplate, Turing.rt_copy_seam, Turing.rt_exists_scratchTemplateTM, Turing.rt_initRecord_fields, Turing.rt_initRestExit, Turing.rt_initRestState, Turing.rt_initRestTM, Turing.rt_initRest_run, Turing.rt_template_fields, Turing.seamCompTM}
\difficulty{7}
For every natural number $k$, finite types $Q$ and $T$ with decidable equality, initial
state $q_0\in Q$, and tag $\tau_0\in T$, there are a natural number $K$, a finite state
type $P$ with decidable equality, a deterministic Turing machine $R$ with $K+2$ work
tapes over the record alphabet, entry and exit states, and a natural number $C$ such
that the following holds. For every binary input word $x$, natural number $s$, and
binary word $w$ with $\operatorname{ofDigits}_2(w)=s$, if a configuration $d$ on the
input obtained by embedding the digits of $x$ has entry state and input-head position
$1$, tape $1$ agrees with the word $w$ stored from position zero on all cells from one
position left of its head through $|w|$ positions right, tape $0$ is blank from one
position left of its head through the record length
$L=k(2s+1)+\mathrm{posWidth}(|x|)+k+4$ positions right, and every tape of index at least
$2$ is blank everywhere, then for some time $t\le C(L+|w|+s+|x|+1)$ the machine reaches
the configuration obtained by changing $d$ only to enter the exit state and to write on
tape $0$ the joined fields of the initial configuration with state $q_0$, input $x$,
blank work tapes, and tag $\tau_0$, placed over the $L$ cells starting at tape $0$'s
initial head position. The machine is not in the exit state at any earlier time, and at
every time through $t$ each work-tape head is at distance at most $C(L+|w|+s+1)$ from
its starting position.
\end{theorem}

\begin{theorem}[Machine for Testing Record Validity]
\lean{Turing.exists_recordValidTM}
\label{Turing.exists_recordValidTM}
\leanok
\uses{Turing.Cfg, Turing.MultiTapeTM, Turing.MultiTapeTM.runFrom, Turing.RecSym, Turing.RecordAt, Turing.RecordValid, Turing.joinFields, Turing.recordTemplate}
For every natural number $k$ and finite decidable types $Q,T$, there are a natural
number $K$, a finite decidable state type $P$, a deterministic machine $R$ with $K+2$
work tapes over the record alphabet, entry and distinct yes/no states, and a constant
$C$ with the following property. For any natural numbers $n,s$, list of fields $f$ of
length $k+3$, and entry-state configuration $d$ whose tape $0$ encodes $f$, tape $1$
encodes the prescribed template for $n,s$, and all tapes with index at least $2$ are
blank, the machine reaches in at most $C(|\langle f\rangle|+|\langle H\rangle|+n+1)$
steps the original configuration with state changed to yes if $f$ encodes a valid record
for parameters $k,n,s,Q,T$, and to no otherwise. It visits neither decision state
earlier, and through that time every work-tape head moves by at most $C(|\langle
f\rangle|+|\langle H\rangle|+1)$ from its initial position, where $H$ is the prescribed
template and validity means that $f$ is the record of a configuration on an input of
length $n$ whose work-tape heads and nonblank cells lie within distance $s$ of the
origin.
\end{theorem}

\begin{theorem}[Enumeration and Successor of Valid Records]
\lean{Turing.exists_recordNextTM}
\label{Turing.exists_recordNextTM}
\leanok
\uses{Turing.Cfg, Turing.MultiTapeTM, Turing.MultiTapeTM.runFrom, Turing.RecSym, Turing.RecordAt, Turing.RecordValid, Turing.fieldsTape, Turing.joinFields, Turing.recordLength, Turing.recordTemplate}
For every natural number $k$ and finite decidable types $Q,T$, there is a family of
finite lists $E_{n,s}$ of records, with no repetitions, such that a record belongs to
$E_{n,s}$ exactly when it encodes a configuration on a binary input of length $n$,
together with a tag, whose work-tape heads and nonblank cells lie within distance $s$ of
the origin. There are also a natural number $K$, a finite decidable state type $P$, a
deterministic machine $R$ with $K+2$ work tapes over the record alphabet, distinct
decision states $\mathrm{yes},\mathrm{no}$, designated states
$\mathrm{first},\mathrm{next}$, and a constant $C$ that perform the following bounded
transitions. On a configuration in state $\mathrm{first}$ with the prescribed blank
first tape, the machine writes the first record of $E_{n,s}$ on that tape and enters
$\mathrm{yes}$, or enters $\mathrm{no}$ if the list is empty. On a configuration in
state $\mathrm{next}$ whose first tape encodes its $i$th record, it writes the next
record and enters $\mathrm{yes}$, or enters $\mathrm{no}$ when there is no next record.
In both cases it preserves all other configuration data, avoids both decision states
before the final time, and keeps every head within distance $C(L+|\langle H\rangle|+1)$
of its initial position through that time, where $L$ is the prescribed record length and
$H$ is the fixed template record; the time is at most $C^{L+|\langle H\rangle|+1}$. The
input record is encoded on work tape $1$ according to the template, and all tapes of
index at least $2$ are blank.
\end{theorem}
